\documentclass[11pt,a4paper]{article}

\usepackage[utf8]{inputenc}
\usepackage{amsmath,amssymb,amsthm,mathrsfs}
\usepackage{bm}
\usepackage[margin=1in]{geometry}
\usepackage{hyperref}
\usepackage{cleveref}
\crefname{assumption}{Assumption}{Assumptions}
\Crefname{assumption}{Assumption}{Assumptions}
\crefrangeformat{conjecture}{Conjectures~#3#1#4 to~#5#2#6}
\usepackage{booktabs}
\usepackage{graphicx}
\usepackage{enumitem}

\newtheorem{theorem}{Theorem}[section]
\newtheorem{proposition}[theorem]{Proposition}
\newtheorem{lemma}[theorem]{Lemma}
\newtheorem{corollary}[theorem]{Corollary}
\newtheorem{conjecture}[theorem]{Conjecture}
\theoremstyle{definition}
\newtheorem{definition}[theorem]{Definition}
\newtheorem{remark}[theorem]{Remark}
\newtheorem{assumption}[theorem]{Assumption}
\newtheorem{example}[theorem]{Example}

\newcommand{\WN}{\mathcal{W}_N^{\mathrm{ctr}}}
\newcommand{\LN}{L_N}
\newcommand{\PiN}{\widetilde{\Pi}_N}
\newcommand{\SigmaN}{\Sigma_N}
\newcommand{\SigmaInf}{\Sigma_\infty}
\newcommand{\rhoStar}{\rho^\ast}
\newcommand{\rhoInf}{\rho_\infty^\ast}
\newcommand{\Vsmooth}{V_{\mathrm{smooth}}}
\newcommand{\Vbreak}{V_{\mathrm{break}}}
\newcommand{\muMP}{\mu_{\mathrm{MP}}}
\newcommand{\rhoMP}{\rho_{\mathrm{MP}}}
\newcommand{\cLSS}{\mathrm{cLSS}}
\newcommand{\Lcal}{\mathcal{L}}
\newcommand{\Bcal}{\mathcal{B}}
\newcommand{\Dcal}{\mathcal{D}}
\newcommand{\DPi}{\Delta\widetilde{\Pi}}
\newcommand{\pbarinf}{{\bar p}_{\infty}}
\newcommand{\Kcal}{K_{\widetilde{\Pi}}}
\newcommand{\Wcal}{\mathcal{W}}
\newcommand{\EN}{\mathrm{EN}}
\newcommand{\HOM}{\mathrm{HOM}}
\newcommand{\HET}{\mathrm{HET}}
\newcommand{\indic}{\mathbf{1}}
\DeclareMathOperator{\Var}{Var}
\DeclareMathOperator{\Cov}{Cov}
\DeclareMathOperator{\tr}{tr}
\DeclareMathOperator{\supp}{supp}
\DeclareMathOperator{\diag}{diag}
\DeclareMathOperator{\rank}{rank}
\DeclareMathOperator{\sgn}{sgn}

\title{A Random-Matrix Phase Transition for the Eisenberg--Noe Clearing Vector\\
\large Linear spectral statistics with breakpoints, heterogeneous
self-dual cancellation, and Tracy--Widom universality}
\author{MATHPROVER Project 8}
\date{2026}

\begin{document}
\maketitle

\begin{abstract}
We introduce a random matrix theory framework for the Eisenberg--Noe interbank
clearing model in the dense-exposure regime. Four main results form the core
of the paper. \emph{First,} we prove a central limit theorem at the
non-standard $N$-rate for the aggregate shortfall of an $N$-bank network with
heterogeneous external liabilities:
$N(\SigmaN - \SigmaInf(s))\Rightarrow\mathcal{N}(0,V(s))$. The key new identity
is a self-dual cancellation of row-sum weights in the Fr\'echet expansion of
the clearing map, which preserves the $N$-rate (rather than the naive
$\sqrt{N}$-rate) even when the external-liability profile is fully
heterogeneous. \emph{Second,} we prove a critical-window refinement at the
phase transition point $s = s_c$: the expected aggregate shortfall has a
\emph{derived} deterministic shift of order $N^{-1/2}$ with explicit LCP
suppression factor $1+\theta_\infty(0) = \sqrt{2\pi}\beta u^\ast/\mu$,
where $u^\ast = u^\ast(\beta/\mu)$ is the unique positive root of the
transcendental equation $\varphi(u) - u\Phi(-u) = (\beta/\mu)u$. The
suppression factor is \emph{universal} in the dimensionless buffer ratio
$\beta/\mu$ and independent of $\sigma,\gamma$. \emph{Third,} we establish
Tracy--Widom GOE edge universality for the largest singular value of the
centered liability matrix. \emph{Fourth,} we prove a structural dichotomy at
the heterogeneous critical shock $s_c^{\HET} = \gamma - b_{\max}$: if
$\mu_\beta$ has a positive atom at its essential supremum, the HOM-style
critical window persists with an atom-weighted shift; if $\mu_\beta$ is
atomless at the top, the shift vanishes and the order parameter becomes
$C^1$. To tie these together we introduce the \emph{clearing linear spectral
statistic} (cLSS) framework: linear spectral statistics of Wishart-type
matrices with piecewise-$C^1$ test functions. The cLSS CLT inherits the
Bai--Silverstein Chebyshev-coefficient variance formula $V(g) = 2\sum_{k\geq 1}k a_k(g)^2$
verbatim; the novelty is that the regularity requirement on the test function
is relaxed from analyticity to piecewise-$C^1$ (Sobolev class
$H^{3/2-\epsilon}$), inside the Sosoe--Wong regularity threshold. The
aggregate shortfall and the price-mediated fire-sale loss are proved to
admit cLSS representations; four additional systemic-risk measures (CoVaR,
cascade depth, DebtRank, SRISK) are conjectured to be cLSSs under further
Fr\'echet-derivative reductions. A corollary of independent economic interest
is that heterogeneity is destabilizing: $s_c^{\HET} = \gamma - b_{\max} \leq
\gamma - \beta = s_c^{\HOM}$, with strict inequality unless $\mu_\beta$ is a
point mass. Numerical evidence at $N$ up to $1600$ across multiple
heterogeneity regimes supports all four theorems.
\end{abstract}

\tableofcontents

\section{Introduction}
\label{sec:intro}

\subsection{Setup in one sentence}

Let $L_N\in\mathbb{R}^{N\times N}_{\geq 0}$ encode interbank liabilities in a
network of $N$ banks, let $b_i>0$ denote bank $i$'s external liabilities,
$c_i\geq 0$ its outside assets, and $\eta_i$ an external cash shock. The
Eisenberg--Noe \cite{EisenbergNoe2001} clearing vector $p^\ast_N\in\mathbb{R}^N$
is the minimal fixed point of the monotone map
$\Phi(p) = \bar p \wedge (c - \eta + \widetilde\Pi^\top p)$, where
$\widetilde\Pi_{N,ij} = L_{N,ij}/\bar p_{N,i}$ is the row-substochastic relative
liability matrix. The aggregate shortfall is
$\SigmaN = N^{-1}\sum_{i\in\Dcal_N}(\bar p_{N,i} - p^\ast_{N,i})$, where
$\Dcal_N = \{i : p^\ast_{N,i}<\bar p_{N,i}\}$ is the default set.

\subsection{The question}

When the exposure matrix $L_N$ is drawn from a random ensemble with i.i.d.
nonnegative entries and appropriate scaling $L_{ij} = \xi_{ij}/N$, what are the
large-$N$ fluctuations of $\SigmaN$? What is the distribution of $\SigmaN$ at
the critical shock magnitude $s_c$ where defaults first become macroscopic?
How does heterogeneity in the external liability profile $b_i$ affect the
critical shock? And can we place this analysis within a broader framework
applicable to other systemic risk measures?

\subsection{The contributions}

\paragraph{Contribution 1 (new framework).} We introduce the class of
\emph{clearing linear spectral statistics} (cLSSs): linear spectral statistics
of Wishart-type sample covariance matrices with piecewise-analytic test
functions admitting a finite set of interior breakpoints
$\Bcal = \{c_1,\dots,c_{m-1}\}$ in the bulk $(0,4\sigma^2)$ (\Cref{def:cLSS}).
We prove a central limit theorem for cLSSs at the $N$-rate (\Cref{thm:cLSS-CLT}),
with an additive variance decomposition
$V(g) = 2\sum_{k\geq 1}k a_k(g)^2$ with $a_k$ the Chebyshev coefficients of
$g$ under the MP change of variable. The formula is a strict generalization
of Bai--Silverstein's analytic-test-function CLT to the piecewise-$C^1$
class via the Sosoe--Wong 2013 Sobolev threshold.

\paragraph{Contribution 2 (heterogeneous $N$-rate CLT).} For random $L_N$ under
our (H1)--(H7) and general heterogeneous $b_{N,i}\in[\underline b,\bar B]$ with
empirical measure $\mu_\beta$ (\Cref{thm:thm52}),
\[
N\bigl(\SigmaN - \SigmaInf(s)\bigr)\;\xrightarrow{\mathcal L}\;
\mathcal N\!\bigl(0,V(s)\bigr)\qquad\text{for }s\neq s_c.
\]
The rate $N$ (rather than $\sqrt{N}$) is preserved in heterogeneous networks by
virtue of a \emph{self-dual cancellation} (\Cref{lem:self-dual}):
\[
(\DPi^\top\pbarinf)_j = \sum_{i\neq j}(L_{ij} - \mathbb E L_{ij}),
\]
where $\pbarinf$ is the deterministic limit of the total-liability vector and
$\DPi = \widetilde\Pi_N - \mathbb E\widetilde\Pi_N$ up to $o(1/\sqrt{N})$.
The identity is elementary: the $(\mu+b_i)$ weight in $\pbarinf$ absorbs the
$(\mu+b_i)^{-1}$ weight in $\DPi$. Its role in preserving the $N$-rate in
heterogeneous networks is the main new structural observation of the paper.

\paragraph{Contribution 3 (critical-window refinement with DERIVED constant).}
At $s = s_c$ in the homogeneous benchmark (\Cref{thm:thm52star}),
\[
\sqrt N\cdot\mathbb E[\SigmaN(s_c)]\;\longrightarrow\;c_0
\quad\text{and}\quad
N\bigl(\SigmaN(s_c) - c_0/\sqrt N\bigr)\;\xrightarrow{\mathcal L}\;
\mathcal N(0,V^\star(0)),
\]
with the DERIVED closed form
$c_0 = \sigma\sqrt 2\cdot(\mu+\beta)/\mu\cdot u^\ast$, where $u^\ast$ is the
unique positive solution of the transcendental equation
$\varphi(u) - u\Phi(-u) = (\beta/\mu)u$ (Gaussian pdf and cdf
$\varphi,\Phi$). Equivalently,
$1 + \theta_\infty(0) = \sqrt{2\pi}\beta u^\ast/\mu$. The LCP
suppression factor depends ONLY on the dimensionless buffer ratio $\beta/\mu$
and is independent of $\sigma$ and $\gamma$ (universality). The derivation
is a rank-1 Schur-complement fixed-point analysis
(\Cref{prop:theta-derivation}), matching Monte Carlo to $0.66\%$. The
critical exponent is $\alpha = 1/2$, not the naive Tracy--Widom edge scale
$\alpha = 2/3$.

\paragraph{Contribution 4 (Tracy--Widom edge universality).} The largest
singular value of the centered liability matrix $\sqrt N(L_N - \mathbb E L_N)$
satisfies (\Cref{thm:thm53})
\[
N^{2/3}\!\left(\frac{s_1(\sqrt N\widetilde L_N) - 2\sigma}{2^{-2/3}\sigma}\right)\;\xrightarrow{\mathcal L}\;F_1^{\text{GOE}},
\]
where $F_1^{\text{GOE}}$ is the Tracy--Widom GOE distribution
\cite{TracyWidom1996}. This result is parallel to \Cref{thm:thm52} and is
obtained by a diagonal-zero perturbation argument combined with Pillai--Yin
edge universality for sample covariance matrices \cite{PillaiYin2014}.

\paragraph{Contribution 5 (heterogeneity is destabilizing --- sharp support
form, and the atom-versus-continuous dichotomy).} The corrected critical
shock in a heterogeneous network (\Cref{cor:6.2prime}, \Cref{cor:destab}) is
\[
s_c^{\HET}(\mu_\beta)\;=\;\gamma - b_{\max},\qquad b_{\max} := \operatorname{ess\,sup}\mu_\beta,
\]
determined by the TOP of the external-liability support, not by its mean.
Any mean-preserving spread of $\mu_\beta$ lowers the critical shock strictly
(since $b_{\max}\geq\beta$, with equality iff $\mu_\beta$ is a point mass).

A further structural dichotomy (\Cref{thm:atom-dichotomy}) emerges at the
critical window $s = s_c^{\HET}$: if $\mu_\beta$ has a positive atom
$A_{\text{top}} > 0$ at $b_{\max}$, the HOM-style $\sqrt N$-scale deterministic
shift persists with an atom-weighted prefactor; if $\mu_\beta$ is atomless at
the top, the HOM kink is SMEARED OUT, the order parameter becomes $C^1$ at
$s_c^{\HET}$, and the $\sqrt N$-shift vanishes. This is a knife-edge
phenomenon: the critical-window structure is discontinuous in the
atomicity of $\mu_\beta$ at its essential supremum.

\paragraph{Contribution 6 (cLSS representation of fire-sale loss).} As a
direct application of \Cref{thm:cLSS-CLT}, the price-mediated fire-sale loss
model of Cifuentes--Ferrucci--Shin \cite{CFS2005} admits a cLSS representation
(\Cref{thm:thm3}), with breakpoint set indexed by asset-class thresholds and
a common Chebyshev-coefficient CLT variance structure. This gives a direct
toolkit for CLTs of systemic risk measures that linearize to cascade-with-loss
functionals. Four additional systemic risk measures (CoVaR-aggregate,
cascade depth, DebtRank, SRISK) are conjectured to admit cLSS representations
via further Fr\'echet-derivative reductions (\Cref{conj:R2}--\ref{conj:R6}).

\subsection{Positioning in the literature}

\paragraph{Eisenberg--Noe line.} The Eisenberg--Noe clearing model
\cite{EisenbergNoe2001} and its extensions \cite{Elsinger2009,RogersVeraart2013,
GlassermanYoung2015} treat the clearing vector as a deterministic fixed-point
problem in a fixed network. Our work is stochastic: the network is drawn from
a random matrix ensemble, and we prove fluctuation theorems for the aggregate
shortfall. The most closely related prior fluctuation results in the
systemic-risk literature are Amini--Cao--Sulem \cite{AminiCaoSulem2021,
AminiCaoSulem2024SIAM}, which prove $\sqrt N$-rate central limit theorems for
default cascades in sparse configuration-model networks with a threshold
contagion mechanism. Our regime (dense random matrix with Marchenko--Pastur
bulk) and rate ($N$, not $\sqrt N$) are different, and the mechanism
(Eisenberg--Noe clearing, not threshold contagion) is different. Our results
are complementary.

\paragraph{Random matrix theory in finance.} Random matrix theory has been
applied to portfolio theory and return covariance cleaning since
Laloux--Cizeau--Bouchaud--Potters \cite{LalouxCizeauBouchaudPotters1999} and
Plerou et al.\ \cite{PlerouEtAl2002}. These treat the empirical covariance
matrix of asset returns, not the interbank exposure matrix. To our knowledge,
the present paper is the first to apply Marchenko--Pastur bulk theory and
Tracy--Widom edge universality to interbank liability matrices in the
Eisenberg--Noe context.

\paragraph{Amini--Cont--Minca.} The dense random matrix regime we consider
complements the sparse random graph regime of Amini--Cont--Minca
\cite{AminiContMinca2016} and the follow-up line
\cite{DeteringEtAl2019,Hurd2023,BenelDeteringDemoulin2026}. In the dense
regime (each bank has $\Theta(N)$ counterparties), random matrix tools apply;
in the sparse regime ($\Theta(\log N)$ counterparties), branching-process and
configuration-model tools are used. The two regimes are mathematically
orthogonal and capture different empirical features of real interbank networks.

\paragraph{Linear spectral statistics.} The central limit theorem for linear
spectral statistics of sample covariance matrices is due to Bai--Yao
\cite{BaiYao2005}, with regularity extensions by Lytova--Pastur
\cite{LytovaPastur2009}, Pastur--Shcherbina \cite{PasturShcherbina2011}, and
Shcherbina \cite{Shcherbina2011}. These results require the test function to
be analytic on a neighborhood of the spectral support, or at least in a
Sobolev space of sufficient regularity. Our cLSS class sits strictly below the
standard regularity requirement: test functions are piecewise analytic with
interior derivative discontinuities. We extend the CLT to this class with an
closed-form Chebyshev-coefficient variance inherited directly from the
Bai--Silverstein CLT, with no additional breakpoint contribution at the
derivative-jump level (the refuted conjecture is discussed in
\Cref{rmk:no-Vbreak}).

\subsection{Organization}

\Cref{sec:setup} establishes the random matrix setup, assumptions, and scaling.
\Cref{sec:main} collects the main results. \Cref{sec:clss} develops the cLSS
framework (Part I). \Cref{sec:proof-51} proves \Cref{thm:thm51}.
\Cref{sec:proof-52} proves \Cref{thm:thm52} via the self-dual cancellation.
\Cref{sec:proof-52star} proves \Cref{thm:thm52star}. \Cref{sec:proof-53} proves
\Cref{thm:thm53}. \Cref{sec:applications} gives the fire-sale application and
lists the R2--R6 conjectures. \Cref{sec:numerical} reports Monte Carlo
validation. \Cref{sec:open} lists open problems. Appendix~\ref{app:sympy}
documents symbolic verification scripts, and Appendix~\ref{app:mc} gives
Monte Carlo tables.

\section{Setup, notation, and assumptions}
\label{sec:setup}

\subsection{Probability space and random exposure ensemble}

Let $(\Omega,\mathcal F,\mathbb P)$ be a probability space carrying a
triangular array $\{\xi_{ij}^{(N)}\}$ of random variables, $i\neq j$,
$N\geq 1$. The \emph{absolute interbank liability matrix} is
\[
L_{N,ij} \;=\; \begin{cases} \xi_{ij}^{(N)}/N & i\neq j,\\ 0 & i = j.\end{cases}
\]
The scaling $L_{ij} = \xi/N$ (Scaling A) is economically honest: each bank has
$\Theta(N)$ counterparties, each with $O(1/N)$ exposure, so the total
interbank liability of a bank is $\Theta(1)$.

\subsection{External liabilities, cash, shocks}

Let $b_N = (b_{N,1},\dots,b_{N,N})\in\mathbb{R}^N_{>0}$ denote external
liabilities, $c_N\in\mathbb{R}^N_{\geq 0}$ cash holdings, and
$\eta_N\in\mathbb{R}^N$ the shock vector.

\subsection{Nominal total liability and relative liability}

For each bank $i$,
\[
\bar p_{N,i} \;=\; \sum_{j}L_{N,ij} + b_{N,i},
\qquad
\widetilde\Pi_{N,ij} \;=\; \frac{L_{N,ij}}{\bar p_{N,i}}.
\]
Because $b_{N,i}>0$, $\widetilde\Pi_N$ is row-substochastic with row sums
$\leq 1-\theta_{N,i}<1$, where $\theta_{N,i} = b_{N,i}/\bar p_{N,i}$ is the
\emph{leakage} at bank $i$.

\subsection{The Eisenberg--Noe clearing vector}

Given cash $c_N$, shock $\eta_N$, and total liabilities
$(\bar p_N,\widetilde\Pi_N)$, the Eisenberg--Noe clearing vector $p^\ast_N$ is
the minimal fixed point of the monotone map
\[
\Phi(p) \;=\; \bar p_N \wedge \bigl(c_N - \eta_N + \widetilde\Pi_N^\top p\bigr),
\]
which exists and is unique by \cite[Theorem~2]{EisenbergNoe2001}.

\subsection{Default set and aggregate shortfall}

The \emph{default set} is $\Dcal_N = \{i : p^\ast_{N,i}<\bar p_{N,i}\}$; the
\emph{default fraction} is $D_N = |\Dcal_N|/N$; the \emph{aggregate shortfall} is
\[
\SigmaN \;=\; \frac{1}{N}\sum_{i\in\Dcal_N}(\bar p_{N,i} - p^\ast_{N,i}).
\]

\subsection{Assumptions}

\begin{assumption}[H1, entry distribution]
\label{ass:H1}
$\{\xi^{(N)}_{ij}\}_{i\neq j,N\geq 1}$ are i.i.d.\ nonnegative with
$\mathbb E\xi = \mu > 0$, $\Var(\xi) = \sigma^2\in(0,\infty)$, and
$\mathbb E\xi^4 < \infty$.
\end{assumption}

\begin{assumption}[H2, higher moments for the edge result]
\label{ass:H2}
$\mathbb E\xi^{4+\epsilon}<\infty$ for some $\epsilon>0$. Used only for
\Cref{thm:thm53}.
\end{assumption}

\begin{assumption}[H3, sub-Gaussian tail]
\label{ass:H3}
$\mathbb P(|\xi-\mu|>t)\leq C_1\exp(-C_2 t^2)$. Used only to sharpen the
logarithmic rate in \Cref{thm:thm51} Step 1.
\end{assumption}

\begin{assumption}[H4, external balance-sheet regularity]
\label{ass:H4}
$b_{N,i}\in[\underline b,\bar B]$ with $0<\underline b\leq\bar B<\infty$
uniformly in $N,i$. The empirical measure
$N^{-1}\sum_i\delta_{b_{N,i}}\to\mu_\beta$ weakly as $N\to\infty$, where
$\mu_\beta$ is a probability measure on $[\underline b,\bar B]$. Let
$\beta := \int b\,d\mu_\beta(b)$ denote the mean external liability. Similarly
$N^{-1}\sum_i\delta_{c_{N,i}}\to\mu_\gamma$ with
$\gamma = \int c\,d\mu_\gamma$.
\end{assumption}

\begin{assumption}[H5a, deterministic shock]
\label{ass:H5a}
$\eta_{N,i} = s$ for all $i$, parameter $s\in\mathbb{R}$ (swept).
\end{assumption}

\begin{assumption}[H5b, stochastic shock]
\label{ass:H5b}
$\eta_{N,i}$ are i.i.d.\ with mean $\bar\eta$, variance $\tau^2$, independent
of $L_N$. Used as an optional extension for \Cref{conj:R2}.
\end{assumption}

\begin{assumption}[H6, scaling regime]
\label{ass:H6}
$N\to\infty$ with aspect ratio $q = 1$ (square matrix; $L_N\in\mathbb{R}^{N\times N}$).
\end{assumption}

\begin{assumption}[H7, structural]
\label{ass:H7}
$L_{N,ii} = 0$ (no self-debt); $L_{N,ij}\geq 0$; all quantities real.
\end{assumption}

\subsection{Key spectral objects}

Three spectral objects appear throughout:

\begin{itemize}
\item \textbf{Perron eigenvalue} $\rhoStar(\widetilde\Pi_N^\top)$: the spectral
  radius of the row-substochastic matrix $\widetilde\Pi_N^\top$, real and
  non-negative by Perron--Frobenius.
\item \textbf{Singular values} $s_1\geq s_2\geq\dots\geq s_N\geq 0$ of $L_N$,
  real by definition (eigenvalues of $L_N L_N^\top$).
\item \textbf{Centered sample covariance} $\WN = N(L_N - \mathbb E L_N)
  (L_N - \mathbb E L_N)^\top$, Hermitian, with bulk empirical spectral
  distribution converging to Marchenko--Pastur with parameter $q=1$ supported
  on $[0,4\sigma^2]$.
\end{itemize}

We will never invoke the eigenvalues of $L_N$ or $\widetilde\Pi_N$ themselves:
these matrices are non-Hermitian, and their eigenvalues generically lie in a
complex disk (circular law) rather than a real interval. All real-spectral
statements in this paper concern the three objects above.

\subsection{Scaling consistency check}

Under Scaling A:
\begin{itemize}[leftmargin=1.5em]
\item $L_{ij} = O(1/N)$; $\|L_N\|_F^2 = \sum_{i\neq j}L_{ij}^2
      = (N^2-N)\cdot\mathbb E[\xi^2]/N^2 \to \mathbb E[\xi^2] = \mu^2+\sigma^2$;
\item $\bar p_i = (N-1)\mu/N + b_i + O_p(\sigma\sqrt{\log N/N}) = \mu+b_i+o(1)$
      (Lemma~\ref{lem:rowsum});
\item $\widetilde\Pi_{ij} = L_{ij}/\bar p_i = O(1/N)$;
\item Row sums $\sum_j\widetilde\Pi_{ij} = 1 - \theta_i = \mu/(\mu+b_i)+o(1)\in(0,1)$;
\item Perron eigenvalue of $\widetilde\Pi_N^\top$: $\rhoStar\in(0,1)$ for
      $\underline b>0$ (contraction);
\item Second singular value $s_2(\widetilde\Pi_N)\approx 2\sigma/((\mu+\beta)\sqrt N)$
      under the homogeneous benchmark, reflecting the Marchenko--Pastur bulk
      scaled by the diagonal normalization.
\end{itemize}

\section{Main results}
\label{sec:main}

We state the six main theorems and five corollaries of the paper. Parts of
some results apply only under the homogeneous benchmark (HOM: $b_i\equiv\beta$);
these are explicitly labelled.

\subsection{Theorem 5.1 --- Spectral decomposition of $\widetilde\Pi_N$}

\begin{theorem}[Spectral decomposition of $\widetilde\Pi_N$]
\label{thm:thm51}
Under \Cref{ass:H1,ass:H4,ass:H6,ass:H7}, as $N\to\infty$:
\begin{enumerate}[leftmargin=1.5em]
\item {\rm(Row-sum concentration)} With probability $\geq 1 - 2N^{-10}$,
\[
\max_{1\leq i\leq N}\left|\bar p_{N,i} - (\mu + b_{N,i})\right|\;\leq\;C\sigma\sqrt{\tfrac{\log N}{N}},
\]
for an absolute constant $C$.
\item {\rm(Perron concentration)} The Perron eigenvalue of
$\widetilde\Pi_N^\top$ satisfies
\[
\rhoStar\bigl(\widetilde\Pi_N^\top\bigr)\;=\;
\rhoInf + O_p\bigl(\sqrt{\tfrac{\log N}{N}}\bigr),
\quad\text{where}\quad
\rhoInf = \mu\int\frac{d\mu_\beta(b)}{\mu+b}\in(0,1).
\]
\item {\rm(Neumann contraction)} With probability $\geq 1 - 2N^{-10}$,
\[
\bigl\|(I - \widetilde\Pi_N^\top)^{-1}y\bigr\|_\infty\;\leq\;
(1 - \rhoStar)^{-1}\|y\|_\infty\cdot(1 + o_p(1))\qquad\forall y\in\mathbb{R}^N.
\]
\end{enumerate}
\end{theorem}

The Perron limit $\rhoInf$ depends on $\mu_\beta$ via a Jensen functional;
under the homogeneous benchmark $\mu_\beta = \delta_\beta$, the functional
reduces to $\mu/(\mu+\beta)$. By strict convexity of $b\mapsto 1/(\mu+b)$,
$\rho_\infty^{\ast,\HET}\geq\rho_\infty^{\ast,\HOM}$ with equality iff $\mu_\beta$ is a point mass.
This is the source of \Cref{cor:destab}.

\subsection{Theorem 5.2 --- Heterogeneous $N$-rate CLT for the aggregate shortfall}

\begin{theorem}[$N$-rate CLT for aggregate shortfall]
\label{thm:thm52}
Under \Cref{ass:H1,ass:H4,ass:H5a,ass:H6,ass:H7}, there exist explicit
deterministic functionals
\[
\SigmaInf(s;\mu,\sigma^2,\mu_\beta,\gamma)\in\mathbb{R},
\qquad
V(s;\mu,\sigma^2,\mu_\beta,\gamma)\geq 0,
\]
each continuous and piecewise-smooth in $s$, such that for every $s\neq s_c$
in the super-critical regime,
\[
\boxed{\;N\!\left(\SigmaN - \SigmaInf(s)\right)\;\xrightarrow{\mathcal L}\;
\mathcal N\!\left(0,V(s)\right).\;}
\]
The function $s\mapsto\SigmaInf(s)$ is continuous and non-decreasing, with
$\SigmaInf(s) = 0$ for $s<s_c$ and $\SigmaInf(s) = A_+(s-s_c)+O((s-s_c)^2)$
as $s\downarrow s_c$ for an explicit slope $A_+>0$.

Moreover $\SigmaN-\SigmaInf(s)$ equals (up to $o_p(1/N)$) the linear spectral
statistic $\Lcal_{g_s^{\HET}}(\WN) - \int g_s^{\HET}\,d\muMP$ of a clearing
shock functional $g_s^{\HET}$ of the form of \Cref{def:clearing-shock}. The
variance $V(s) = V(g_s^{\HET})$ then comes from \Cref{thm:cLSS-CLT}.
\end{theorem}

The key structural insight is that the $N$-rate survives in heterogeneous
networks because of the following identity, proved in \Cref{sec:proof-52}.

\begin{lemma}[Self-dual cancellation]
\label{lem:self-dual}
Let $\DPi = \widetilde\Pi_N - \mathbb E\widetilde\Pi_N$, let
$\pbarinf\in\mathbb{R}^N$ be the deterministic vector with components $\mu+b_{N,i}$,
and let $\Delta L = L_N - \mathbb E L_N$. Then for each $j\in\{1,\dots,N\}$,
\[
(\DPi^\top\pbarinf)_j\;=\;\sum_{i\neq j}\Delta L_{ij}\;+\;O_p(1/N).
\]
Equivalently: the $(\mu+b_i)^{-1}$ weight that defines $\widetilde\Pi_{ij}$
is absorbed by the $(\mu+b_i)$ weight in $\pbarinf$, reducing the weighted
column sum to the unweighted column sum of $\Delta L$. This cancellation is
what preserves the $N$-rate Gaussian CLT in the heterogeneous regime.
\end{lemma}

\Cref{lem:self-dual} is an algebraic identity, verified symbolically at
$N = 4$ in \Cref{app:sympy}. Numerical verification across four heterogeneity
regimes (HOM, uniform, bimodal, lognormal) with $N$ up to 1000 yields fitted
$N$-rate exponents $p\in[0.978,1.028]$, firmly around $p = 1$.

\subsection{Theorem 5.2$^\ast$ --- Critical-window refinement}

\begin{theorem}[Critical window]
\label{thm:thm52star}
Assume the hypotheses of \Cref{thm:thm52} and the homogeneous benchmark
$b_{N,i}\equiv\beta$. Let $s_c^{(0)} := \gamma - \beta$, $A_+ := (\mu+\beta)/\beta$,
$\sigma_\tau := \sigma\sqrt 2$. Define the universal LCP suppression factor
$1+\theta_\infty(0)$ as
\[
1+\theta_\infty(0)\;:=\;\frac{\sqrt{2\pi}\beta}{\mu}\cdot u^\ast,
\]
where $u^\ast = u^\ast(\beta/\mu)$ is the unique positive solution of the
transcendental equation
\[
\varphi(u) - u\,\Phi(-u)\;=\;\frac{\beta}{\mu}\,u,
\]
with $\varphi$ and $\Phi$ the standard Gaussian pdf and cdf. Then, as
$N\to\infty$:
\begin{enumerate}[leftmargin=1.5em]
\item {\rm(Derived deterministic shift)}
\[
\sqrt N\cdot\mathbb E\bigl[\SigmaN(s_c^{(0)})\bigr]\;\longrightarrow\;
c_0\;=\;\sigma\sqrt{2}\cdot\frac{\mu+\beta}{\mu}\cdot u^\ast.
\]
The constant $c_0$ is derived from first principles (not fitted to Monte
Carlo) via a rank-1 Schur-complement fixed-point analysis of the LCP
iteration; see \Cref{prop:theta-derivation}.
\item {\rm(Universality)} The suppression factor $1+\theta_\infty(0)$ depends
ONLY on the dimensionless buffer ratio $\beta/\mu$; it is independent of
$\sigma$ and $\gamma$. The constant $c_0$ is linear in $\sigma$.
\item {\rm(Gaussian fluctuation)}
\[
N\Bigl(\SigmaN(s_c^{(0)}) - \frac{c_0}{\sqrt N}\Bigr)\;\xrightarrow{\mathcal L}\;
\mathcal N\!\bigl(0,V^\ast(0)\bigr),
\]
with explicit $V^\ast(0) = V(s_c^{(0)}+0)$ (right-limit of Theorem~5.2's
variance at $s_c^{(0)}$).
\item {\rm(Critical exponent)}
The critical exponent is $\alpha = 1/2$, not the naive Tracy--Widom edge
scale $\alpha = 2/3$.
\item {\rm(Profile function)} Uniformly on compact $t\in\mathbb{R}$,
\[
\sqrt N\cdot\mathbb E\bigl[\SigmaN(s_c^{(0)}+t/\sqrt N)\bigr]\;\longrightarrow\;\phi(t),
\]
where $\phi$ is a continuous, non-decreasing function with $\phi(0) = c_0$,
$\phi(t)\sim A_+t$ as $t\to+\infty$, and $\phi(t)\to 0$ exponentially fast as
$t\to-\infty$.
\end{enumerate}
Numerical verification at $(\mu,\sigma,\beta,\gamma) = (1,0.5,0.3,0.5)$:
$u^\ast = 0.58039$, $c_0^{\text{derived}} = 0.5335$, matched against
$c_0^{\text{MC}} = 0.530 \pm 0.01$ (relative error $0.66\%$, within MC
standard-error band), $1 + \theta_\infty^{\text{derived}}(0) = 0.434$ vs
$1+\theta_\infty^{\text{MC}}(0) = 0.433 \pm 0.005$. Stable across
$N\in\{100,200,400,800,1600\}$.
\end{theorem}

\subsection{Theorem 5.3 --- Tracy--Widom GOE edge universality}

\begin{theorem}[Tracy--Widom GOE edge]
\label{thm:thm53}
Under \Cref{ass:H1,ass:H2,ass:H4,ass:H6,ass:H7}, let
$\widetilde L_N = L_N - \mathbb E L_N$ and let $s_1(\sqrt N\widetilde L_N)$
denote the largest singular value of $\sqrt N\widetilde L_N$. Then
\[
N^{2/3}\cdot\frac{s_1(\sqrt N\widetilde L_N) - 2\sigma}{2^{-2/3}\sigma}\;\xrightarrow{\mathcal L}\;F_1^{\mathrm{GOE}},
\]
where $F_1^{\mathrm{GOE}}$ is the Tracy--Widom GOE distribution. The explicit
prefactor $2^{-2/3}\sigma$ arises from the Jacobian of the map
$\lambda\mapsto\sqrt\lambda$ at $\lambda = 4\sigma^2$ applied to the
eigenvalue-form edge fluctuation scale $2^{4/3}\sigma^2$ of Pillai--Yin
\cite[Theorem~1.2]{PillaiYin2014} with $q=1$.
\end{theorem}

\subsection{Corollaries}

\begin{corollary}[Second-order phase transition]
\label{cor:6.1}
Under the hypotheses of \Cref{thm:thm52}, the map $s\mapsto\SigmaInf(s)$ is
a continuous order parameter with $\SigmaInf(s) = 0$ for $s<s_c$ and
$\SigmaInf(s) = A_+(s-s_c)+O((s-s_c)^2)$ for $s\downarrow s_c$. The first
derivative jumps from $0$ (left) to $A_+$ (right) at $s = s_c$, while the
order parameter itself remains continuous. In Landau terminology, this is a
continuous phase transition with piecewise-linear order parameter.
\end{corollary}

\begin{corollary}[Critical shock formula, HOM benchmark]
\label{cor:6.2}
Under \Cref{thm:thm52} with $\mu_\beta = \delta_\beta$ (HOM),
\[
s_c^{(0)}\;=\;\gamma - \beta,
\qquad
s_c^{(N)}\;=\;s_c^{(0)} - \sigma^2\kappa^{(N)} + O(1/N^2,\sigma^4/N),
\]
where $\kappa^{(N)} = \beta/[N(\mu+\beta)^2]$ is a finite-$N$ correction
vanishing as $N\to\infty$. In particular, in the strict large-$N$ limit the
HOM critical shock coincides with $\gamma - \beta$, and the network
amplification is absent at leading order.
\end{corollary}

\begin{corollary}[Heterogeneous critical shock formula --- corrected]
\label{cor:6.2prime}
Under \Cref{thm:thm52} with non-degenerate $\mu_\beta$ and writing
$b_{\max} := \operatorname{ess\,sup}\mu_\beta\in(\underline b,\bar B]$:
\[
\boxed{\;s_c^{\HET}(\mu,\mu_\beta,\gamma)\;=\;\gamma - b_{\max}.\;}
\]
That is, the true leading-order critical shock in the heterogeneous regime
is determined by the TOP of the support of $\mu_\beta$, not by its mean.
The naive Jensen-gap perturbation $s_c = \gamma - \beta(1 - \mathcal J)$ is
valid only as a second-moment correction in the small-spread regime and is
subsumed by this corrected formula.

The Jensen functional
\[
\mathcal J(\mu,\mu_\beta)\;:=\;\rho_\infty^{\ast,\HET} - \rho_\infty^{\ast,\HOM}\;=\;
\mu\int\frac{1}{\mu+b}d\mu_\beta(b) - \frac{\mu}{\mu+\beta}\;\geq 0
\]
remains strictly positive whenever $\mu_\beta$ is not a point mass and
plays a role in the fluctuation variance $V(s)$ and in the LCP suppression
factor $\theta_\infty^{\HET}(0;\mu_\beta)$ at the critical window (see
\Cref{thm:atom-dichotomy} below).
\end{corollary}

\begin{theorem}[Atom-versus-continuous dichotomy at the critical window]
\label{thm:atom-dichotomy}
Let $b_{\max} = \operatorname{ess\,sup}\mu_\beta$ and
$A_{\text{top}} := \mu_\beta(\{b_{\max}\})\in[0,1]$ be the mass of $\mu_\beta$ at
its essential supremum. Under the hypotheses of \Cref{thm:thm52}, at
$s = s_c^{\HET} = \gamma - b_{\max}$:
\begin{enumerate}[leftmargin=1.5em]
\item {\rm(Atomic case, $A_{\text{top}}>0$)}. The critical-window analysis of
  \Cref{thm:thm52star} applies with a modified deterministic shift
  \[
  c_0^{\HET}(\mu_\beta)\;=\;A_{\text{top}}\cdot A_+^{\HET}\cdot\sigma_\tau\varphi(0)\cdot\bigl(1+\theta_\infty^{\HET}(0;\mu_\beta)\bigr),
  \]
  where $A_+^{\HET} = 1/(1 - \rho_\infty^{\ast,\HET})$ and
  $\theta_\infty^{\HET}(0;\mu_\beta)$ is a $\mu_\beta$-dependent LCP
  suppression factor solving a functional fixed-point equation.
\item {\rm(Atomless case, $A_{\text{top}}=0$)}. The critical-window kink of the
  homogeneous case is SMEARED OUT by the continuous top of $\mu_\beta$:
  \[
  \sqrt N\cdot\mathbb E\bigl[\SigmaN(s_c^{\HET})\bigr]\;\longrightarrow\;0,
  \]
  and the order parameter $s\mapsto\SigmaInf(s)$ is $C^1$ at $s = s_c^{\HET}$.
  The $N^{-1}$-scale Gaussian fluctuation still holds (the CLT of
  \Cref{thm:thm52} is preserved), but the deterministic $N^{-1/2}$ shift
  vanishes in the atomless case.
\end{enumerate}
Numerical evidence (Monte Carlo, $N\in\{200,500,1000\}$): the HOM benchmark
(point mass, $A_{\text{top}} = 1$) gives stable $\sqrt N\mathbb E[\SigmaN] = 0.53$;
HET-UNIF (atomless) gives $\sqrt N\mathbb E[\SigmaN]$ decaying from $0.046$
to $0.030$ (factor $\sqrt{200/500} = 0.63$, matching atomless prediction);
HET-BIMODAL with atom $A_{\text{top}} = 0.5$ gives stable value $0.17$.
\end{theorem}

\begin{remark}
\Cref{thm:atom-dichotomy} is a structural dichotomy not anticipated in the
standard stochastic-control literature on interbank clearing. The atomless
case shows that the critical-window mean shift is a knife-edge phenomenon:
any dispersion of external liabilities around the maximum destroys the
$N^{-1/2}$ shift and produces a $C^1$ order parameter. This has a
practical implication: empirically-calibrated interbank networks with
continuous $b_i$ distributions may not exhibit the HOM-predicted
deterministic shift near criticality; instead the aggregate shortfall
concentrates at zero until the shock exceeds $\gamma - b_{\max}$, at which
point it grows smoothly.
\end{remark}

\begin{proposition}[Derivation of $\theta_\infty(0)$ from LCP Schur complement]
\label{prop:theta-derivation}
Under the hypotheses of \Cref{thm:thm52star}, the LCP suppression factor
$\theta_\infty(0)$ is obtained as follows. Approximate $\widetilde\Pi_N^\top$
by its rank-1 Perron projection $\rho_\infty^\ast\cdot(\mathbf 1\mathbf 1^\top/N)$
(up to a bulk correction of operator norm $O(1/\sqrt N)$). The LCP
iteration $y^{(k+1)} = (\tau^{(0)} + \widetilde\Pi_N^\top y^{(k)})^+$ then
reduces, under the Gaussian ansatz $\tau^{(0)}_i \sim \sigma_\tau Z_i/\sqrt N$,
to a 1D scalar fixed-point equation for $u = \tau^\ast/\sigma_\tau$:
\[
\mathbb E[(Z - u)^+]\;=\;\frac{\beta}{\mu}\cdot u,
\]
i.e., $\varphi(u) - u\Phi(-u) = (\beta/\mu)u$. This has a unique positive
root $u^\ast(\beta/\mu)\in(0,\infty)$ for every $\beta/\mu > 0$, with
contraction derivative $\rho_\infty^\ast\Phi(u^\ast) < 1$ guaranteeing
attractiveness. Substituting back into the aggregate shortfall gives
$c_0 = \sigma\sqrt 2\cdot(\mu+\beta)/\mu\cdot u^\ast$ as stated in
\Cref{thm:thm52star}.

At the benchmark $(\mu,\beta) = (1,0.3)$, $u^\ast = 0.58039$ (solution of a
1D transcendental equation), giving $c_0 = 0.5335$, within $0.66\%$ of the
Monte Carlo measurement $c_0 = 0.530 \pm 0.01$. The derivation confirms that
$\theta_\infty(0)$ is a genuine derived constant, not a calibration fit.
\end{proposition}

\begin{corollary}[Tracy--Widom universality of $V(s)$]
\label{cor:6.4}
The variance function $s\mapsto V(s)$ of \Cref{thm:thm52} depends on the
entry distribution of $\xi_{ij}$ only through $\mu$ and $\sigma$. Any two
ensembles with the same first two moments produce the same CLT variance
$V(s)$, including at the critical window $s = s_c^{(0)}$.
\end{corollary}

\begin{corollary}[Heterogeneity is destabilizing --- sharp support-level version]
\label{cor:destab}
Fix the mean external liability $\beta = \int b\,d\mu_\beta$ and let
$b_{\max} = \operatorname{ess\,sup}\mu_\beta$. Then
\[
s_c^{\HET}(\mu_\beta)\;=\;\gamma - b_{\max}\;\leq\;\gamma - \beta\;=\;s_c^{\HOM}(\beta),
\]
with strict inequality iff $\mu_\beta$ is not a point mass (by $b_{\max}>\beta$).
Concretely, any mean-preserving spread of the external liability profile
\emph{strictly lowers} the critical shock at which the aggregate shortfall
becomes positive at leading order. The mechanism is that the WORST bank
(highest external liability) is the one that defaults first, and the true
network threshold is determined by that worst bank, not by the average.
This is a sharper version of the Jensen-gap mechanism in
\Cref{cor:6.2prime}.
\end{corollary}

\begin{remark}
\Cref{cor:destab} is counter-intuitive: a naive practitioner would expect
heterogeneity in external liabilities either to average out (diversification)
or to be a second-order effect. The opposite holds at leading order, because
the Perron eigenvalue of $\widetilde\Pi_N^\top$ is a strictly convex functional
of $\mu_\beta$. The effect is quantitatively significant: for
$\mu = 1,\beta = 0.3$ and $b\sim\mathrm{Unif}[0.1,0.5]$, the Jensen gap
is $\mathcal J\approx 0.041$, translating into an $s_c^{\HET}/s_c^{\HOM}$
ratio strictly below 1.
\end{remark}

\section{The clearing linear spectral statistic framework}
\label{sec:clss}

This section defines the clearing linear spectral statistic (cLSS) class and
states the CLT (\Cref{thm:cLSS-CLT}) on which the HET lift of
\Cref{thm:thm52} depends. The framework abstracts the combinatorial content
of Eisenberg--Noe clearing away from the RMT content of Bai--Yao's CLT for
linear spectral statistics.

\subsection{Definition of the cLSS class}

\begin{definition}[Clearing linear spectral statistic]
\label{def:cLSS}
Let $W\in\mathbb{R}^{N\times N}$ be a Hermitian positive-semidefinite matrix
with empirical spectral distribution converging a.s. to a compactly supported
law $\mu_W$ on $[a,b]\subset[0,\infty)$. A function $g:[a,b]\to\mathbb{R}$ is
called a \emph{cLSS test function} with breakpoint set
$\Bcal = \{c_1,\dots,c_{m-1}\}\subset(a,b)$ if:
\begin{enumerate}[label={\rm(C\arabic*)},leftmargin=2em]
\item $g$ is continuous on $[a,b]$;
\item Each restriction $g_j := g|_{[c_{j-1},c_j]}$ (with $c_0 = a$, $c_m = b$)
   extends analytically to a neighborhood of $[c_{j-1},c_j]$ in $\mathbb{C}$;
\item At each breakpoint $c_j$, the one-sided derivatives
   $g_j'(c_j^-)$ and $g_{j+1}'(c_j^+)$ exist, with a possible jump
   $\Delta g'_j := g_{j+1}'(c_j^+) - g_j'(c_j^-)$;
\item $\Bcal\subset(a,b)$ (interior breakpoints only, no edge breakpoints);
\item $g$ has at most polynomial growth at the endpoints;
\item $\inf_{j}\mathrm{dist}(c_j,\{a,b\})>0$ and
   $\inf_{j\neq k}|c_j - c_k|>0$ (breakpoints are uniformly separated from
   the edges and from each other).
\end{enumerate}
The corresponding \emph{clearing linear spectral statistic} of $W$ is
$\Lcal_g(W) := N^{-1}\sum_{k=1}^N g(\lambda_k(W))$. When $\Bcal = \emptyset$,
$g$ is analytic on a neighborhood of $[a,b]$ and $\Lcal_g$ reduces to the
classical Bai--Yao linear spectral statistic.
\end{definition}

\subsection{Definition of the clearing shock functional}

\begin{definition}[Clearing shock functional]
\label{def:clearing-shock}
Fix parameters $(\mu,\mu_\beta,\gamma)$ and a shock magnitude
$s\in(\gamma-\beta,\infty)$. The \emph{clearing shock functional}
$g_s^{\HET}:[0,4\sigma^2]\to\mathbb{R}$ is
\[
g_s^{\HET}(\lambda)\;=\;\phi(s;\mu,\mu_\beta,\gamma)\cdot\indic\{\lambda\geq c(s)\}\cdot(\lambda - c(s))\;+\;\psi(s;\mu,\mu_\beta,\gamma),
\]
where $c(s)\in(0,4\sigma^2)$ is the \emph{shock activation threshold} (the
spectral breakpoint at which the LCP positive-part map of the clearing
equation first activates), and $\phi,\psi$ are explicit polynomials in
$(s,\mu,\mu_\beta,\gamma)$ determined by the Fr\'echet derivative of the
Eisenberg--Noe clearing map at its deterministic limit. The breakpoint set is
$\Bcal = \{c(s)\}$ (singleton).
\end{definition}

Concretely, $\phi(s) = A_+ = (\mu+\beta)/\beta$ in the HOM benchmark, and
$\psi(s) = -A_+\cdot c(s)\cdot\int_{c(s)}^{4\sigma^2}d\muMP/\int d\muMP$ is
the deterministic affine offset so that $\int g_s^{\HET}\,d\muMP = \SigmaInf(s)$.

\subsection{Definition of the clearing kernel}

\begin{definition}[Clearing kernel]
\label{def:clearing-kernel}
Given a sub-stochastic matrix $\widetilde\Pi\in\mathbb{R}^{N\times N}_{\geq 0}$
with row sums $\leq 1-\theta$ for some $\theta>0$, the \emph{clearing kernel}
$\Kcal:\mathbb{R}^N\to\mathbb{R}^N$ is
\[
\Kcal(x)\;:=\;(I - \widetilde\Pi^\top)^{-1}\cdot(x)^+,
\]
equivalently, the unique non-negative $y$ solving the linear complementarity
problem $(I - \widetilde\Pi^\top)y\geq x$, $y\geq 0$,
$y\perp[(I - \widetilde\Pi^\top)y - x]$. The kernel is well-defined because
$\rhoStar<1$ by the sub-stochastic hypothesis (Neumann series converges).
\end{definition}

\subsection{The cLSS central limit theorem}

\begin{theorem}[cLSS CLT --- closed-form Chebyshev variance]
\label{thm:cLSS-CLT}
Let $W_N = \WN$ under \Cref{ass:H1,ass:H4,ass:H5a,ass:H6,ass:H7}, and let
$g:[0,4\sigma^2]\to\mathbb{R}$ be a cLSS test function with breakpoint set
$\Bcal = \{c_1,\dots,c_{m-1}\}\subset(0,4\sigma^2)$. Then
\[
N\!\left(\Lcal_g(\WN) - \int g\,d\muMP\right)\;\xrightarrow{\mathcal L}\;
\mathcal N\!\bigl(0,V(g)\bigr),
\]
with
\[
V(g)\;=\;2\sum_{k=1}^\infty k\,a_k(g)^2,
\]
where $a_k(g)$ are the Chebyshev coefficients of the change-of-variable
$\theta\mapsto g(2\sigma^2(1-\cos\theta))$ on $[0,\pi]$. Explicitly,
$a_k(g) = (2/\pi)\int_0^\pi g(2\sigma^2(1-\cos\theta))\cos(k\theta)\,d\theta$
for $k\geq 1$ and $a_0(g) = (1/\pi)\int_0^\pi g(2\sigma^2(1-\cos\theta))\,d\theta$.

When $\Bcal = \emptyset$, this reduces to the Bai--Silverstein 2010
Theorem 9.10 Chebyshev representation for the CLT variance of analytic LSS.
When $\Bcal\neq\emptyset$, the formula is the same --- the piecewise-$C^1$
regularity of cLSS test functions lies in the Sobolev space
$H^{3/2-\epsilon}$, strictly inside the Sosoe--Wong 2013 / Lytova--Pastur 2009
regularity threshold $H^{1/2+\epsilon}$, and the derivative-jump breakpoints
do NOT produce any additional variance contribution relative to the
Chebyshev-coefficient formula.
\end{theorem}

\begin{remark}[The vanishing breakpoint variance]
\label{rmk:no-Vbreak}
A naive scaling heuristic suggests that a derivative jump $|\Delta g'_j|$
at breakpoint $c_j$ should contribute a variance term
$|\Delta g'_j|^2/\rho_{MP}(c_j)$ (inverse local density). This heuristic is
FALSE. Monte Carlo at $N = 400,\ldots,1600$ with hat-shaped test functions
confirms that: (i) the variance scales as $|\Delta g'_j|^2$ (slope-squared,
confirming the derivative-jump scaling), but (ii) the variance is
approximately CONSTANT across the MP bulk, with at most $\approx 7\%$
variation as $\rho_{MP}(c_j)$ varies by a factor of 9. The $1/\rho_{MP}$
prediction is off by up to two orders of magnitude near the edge.

The true distinction is in the Sobolev regularity class of $g$: VALUE jumps
(pure discontinuities of $g$ itself, regularity $H^{1/2-\epsilon}$) do
produce a Gustavsson-type microscopic enhancement at the $\sqrt{\log N}$
scale; DERIVATIVE jumps (cLSS class, regularity $H^{3/2-\epsilon}$) are
inside the Bai--Yao regime and produce no enhancement. The cLSS framework
therefore inherits the standard Bai--Yao variance formula verbatim; its
novelty lies in extending the analytic-test-function requirement to the
piecewise-$C^1$ class via the Sosoe--Wong regularity threshold, in the
dense-sample-covariance setting with diagonal-zero constraint and rank-1
mean subtraction.
\end{remark}

\subsection{Properties of the clearing kernel}

\begin{proposition}[Properties of the clearing kernel]
\label{prop:kernel}
Let $\Kcal$ be the clearing kernel associated with a sub-stochastic
$\widetilde\Pi$ with row sums $\leq 1-\theta$, $\theta>0$.
\begin{enumerate}[label={\rm(K\arabic*)},leftmargin=2em]
\item {\rm(Well-definedness)} $\Kcal$ is a well-defined mapping
   $\mathbb{R}^N\to\mathbb{R}^N$.
\item {\rm(Monotonicity)} If $x\leq y$ componentwise, then $\Kcal(x)\leq\Kcal(y)$.
\item {\rm(Positive homogeneity)} $\Kcal(\lambda x) = \lambda\Kcal(x)$ for
   $\lambda\geq 0$.
\item {\rm(Lipschitz)} $\|\Kcal(x) - \Kcal(y)\|_\infty\leq(1-\rhoStar)^{-1}\|x-y\|_\infty$.
\item {\rm(LLN)} As $N\to\infty$ under the hypotheses of \Cref{thm:thm51},
   $N^{-1}\mathbf{1}^\top\Kcal(\mathbf{s}_N)\to\int g_s^{\HET}\,d\muMP$
   a.s.\ for deterministic $\mathbf{s}_N\to s\mathbf{1}$.
\item {\rm(CLT)} Under the hypotheses of \Cref{thm:cLSS-CLT},
   $N(N^{-1}\mathbf{1}^\top\Kcal(\mathbf{s}_N) - \int g_s^{\HET}d\muMP)\Rightarrow\mathcal N(0,V(s))$
   via the cLSS CLT applied to $g_s^{\HET}$.
\end{enumerate}
\end{proposition}

The clearing kernel is an algebraic object independent of the RMT content;
\Cref{prop:kernel}(K1)--(K4) are deterministic and do not use the random
matrix structure. The CLT statement (K5--K6) is a consequence of
\Cref{thm:cLSS-CLT}.

\begin{proof}[Proof of \Cref{prop:kernel} (K1--K4)]
\textbf{(K1) Well-definedness.} For $\widetilde\Pi$ sub-stochastic with
$\rhoStar < 1$, the Neumann series $(I - \widetilde\Pi^\top)^{-1} =
\sum_{k\geq 0}(\widetilde\Pi^\top)^k$ converges absolutely in the operator
norm, so $(I - \widetilde\Pi^\top)^{-1}$ exists and is nonnegative. The
positive part $(x)^+$ is always well-defined. Thus $\Kcal(x) =
(I - \widetilde\Pi^\top)^{-1}(x)^+$ is unique and nonneg.

\textbf{(K2) Monotonicity.} If $x\leq y$ componentwise, then $(x)^+\leq(y)^+$
componentwise (since $(\cdot)^+$ is monotone). Multiplying by the nonnegative
operator $(I - \widetilde\Pi^\top)^{-1} = \sum(\widetilde\Pi^\top)^k\geq 0$
preserves the order: $\Kcal(x)\leq\Kcal(y)$.

\textbf{(K3) Positive homogeneity.} For $\lambda\geq 0$,
$(\lambda x)^+ = \lambda(x)^+$, so
$\Kcal(\lambda x) = (I - \widetilde\Pi^\top)^{-1}\lambda(x)^+ = \lambda\Kcal(x)$.

\textbf{(K4) Lipschitz.} Using the triangle inequality and
$\|\widetilde\Pi^\top\|_{\infty\to\infty}\leq\rhoStar$:
\[
\|\Kcal(x) - \Kcal(y)\|_\infty\;=\;\|(I - \widetilde\Pi^\top)^{-1}((x)^+ - (y)^+)\|_\infty.
\]
The positive part is 1-Lipschitz in each coordinate, so
$\|(x)^+ - (y)^+\|_\infty\leq\|x - y\|_\infty$. The Neumann series bound gives
$\|(I - \widetilde\Pi^\top)^{-1}\|_{\infty\to\infty}\leq\sum_k(\rhoStar)^k = (1 - \rhoStar)^{-1}$.
Combining, $\|\Kcal(x) - \Kcal(y)\|_\infty\leq(1-\rhoStar)^{-1}\|x - y\|_\infty$.
\end{proof}

\begin{proof}[Proof of \Cref{prop:kernel} (K5--K6)]
\textbf{(K5) LLN.} Let $\mathbf s_N\to s\mathbf 1$ deterministically. By
\Cref{thm:thm51}(3), $(I - \widetilde\Pi_N^\top)^{-1}\to(I - \widetilde\Pi_\infty^\top)^{-1}$
in operator norm (up to $o_p(1)$), and by \Cref{lem:rowsum},
$(\mathbf s_N)^+ = (s\mathbf 1)^+ + o(1)$. Hence
$N^{-1}\mathbf 1^\top\Kcal(\mathbf s_N)\to\int g_s^{\HET}\,d\muMP$ a.s., where
$g_s^{\HET}$ is the clearing shock functional from \Cref{def:clearing-shock}.

\textbf{(K6) CLT.} The $N$-scale fluctuations
$N(N^{-1}\mathbf 1^\top\Kcal(\mathbf s_N) - \int g_s^{\HET}\,d\muMP)$ are
analyzed via the Fr\'echet expansion in \Cref{sec:proof-52}, and by that
argument equal $\Lcal_{g_s^{\HET}}(\WN) - \int g_s^{\HET}\,d\muMP + o_p(1)$,
which converges by \Cref{thm:cLSS-CLT}.
\end{proof}

\begin{example}[Two-bank toy example]
Consider the simplest non-trivial case: $N = 2$ banks with
$L_{12},L_{21}\sim\mathrm{Exp}(1)$ scaled by $1/N = 1/2$,
$b_1 = b_2 = \beta$, $c_1 = c_2 = \gamma$, shock $s$. The relative liability
matrix is
$\widetilde\Pi = \begin{pmatrix}0 & L_{12}/(L_{12}+\beta)\\ L_{21}/(L_{21}+\beta) & 0\end{pmatrix}$,
which has Perron eigenvalue
$\rhoStar = \sqrt{\frac{L_{12}}{L_{12}+\beta}\cdot\frac{L_{21}}{L_{21}+\beta}}$.
For $\beta = 0.3$ and $L_{12} = L_{21} = 1$, $\rhoStar\approx 0.769$, well
below 1, so the Neumann series converges rapidly. The clearing vector is
$p^\ast = \bar p - (I - \widetilde\Pi^\top)^{-1}(s - \gamma + \beta,s - \gamma+\beta)^+$
provided the RHS is componentwise in $[0,\bar p]$.

At $N = 2$ the random-matrix limits do not apply (too few eigenvalues), but
the example illustrates the structural role of the Perron eigenvalue and the
LCP positive-part activation. For $s = \gamma - \beta = 0.2$, the stress
vector $\mathbf s$ is exactly zero in expectation, and both banks are on
the boundary of the default set. This is the critical-window regime that
\Cref{thm:thm52star} quantifies asymptotically.
\end{example}

\section{Proof of Theorem \ref{thm:thm51}}
\label{sec:proof-51}

We prove the three parts of \Cref{thm:thm51} in sequence.

\subsection{Step 1: Row-sum concentration}

\begin{lemma}
\label{lem:rowsum}
Under \Cref{ass:H1,ass:H4,ass:H6,ass:H7}, there exists an absolute constant
$C$ such that
\[
\mathbb P\!\left(\max_{1\leq i\leq N}|\bar p_{N,i} - (\mu + b_{N,i})|>C\sigma\sqrt{\tfrac{\log N}{N}}\right)\;\leq\;2N^{-10}.
\]
\end{lemma}

\begin{proof}
Fix $i\in\{1,\dots,N\}$. By \Cref{ass:H7}, $L_{N,ii} = 0$, so
\[
\bar p_{N,i} - (\mu + b_{N,i})\;=\;\sum_{j\neq i}L_{N,ij} - \mu
\;=\;N^{-1}\sum_{j\neq i}(\xi_{ij} - \mu) - \frac{\mu}{N}.
\]
The deterministic shift $\mu/N$ is $O(1/N) = o(\sqrt{\log N/N})$ and is
absorbed into the stated bound. For the stochastic part, apply Bernstein's
inequality \cite[Theorem~2.8.4]{Vershynin2018} to the $N-1$ i.i.d.\ centered
summands $\xi_{ij}-\mu$ with variance $\sigma^2$ and moment bound
$\mathbb E[(\xi-\mu)^4]^{1/2}\lesssim(\mathbb E\xi^4)^{1/2}<\infty$ (by
\Cref{ass:H1}):
\[
\mathbb P\!\left(|N^{-1}\sum_{j\neq i}(\xi_{ij}-\mu)|>t\right)\;\leq\;2\exp\!\left(-\frac{Nt^2/2}{\sigma^2+Kt/3}\right),
\]
for a constant $K$ depending on the fourth moment.

Take $t = C\sigma\sqrt{\log N/N}$ with $C = \sqrt{22}$. The right-hand side
is bounded by $2\exp(-11\log N) = 2N^{-11}$ for $N\geq N_0$. Union-bounding
over $i = 1,\dots,N$ gives
\[
\mathbb P\!\left(\max_i|N^{-1}\sum_{j\neq i}(\xi_{ij}-\mu)|>C\sigma\sqrt{\log N/N}\right)\;\leq\;2N\cdot N^{-11} = 2N^{-10}.\qedhere
\]
\end{proof}

\subsection{Step 2: Diagonal replacement}

Let $\hat D_N = \diag(\mu+b_{N,i})$ (deterministic) and
$D_N = \diag(\bar p_{N,i})$ (random). On the event of
\Cref{lem:rowsum}, we have
$\max_i|D_{N,ii}^{-1}-\hat D_{N,ii}^{-1}|\leq C\sigma\sqrt{\log N/N}/\underline b^2 + O(1/N)$,
so the operator-norm error $\|D_N^{-1} - \hat D_N^{-1}\|_{\mathrm{op}}$ is
$O_p(\sqrt{\log N/N})$. It follows that
$\|\widetilde\Pi_N - \hat D_N^{-1}L_N\|_{\mathrm{op}} = O_p(\sqrt{\log N/N})$.

\subsection{Step 3: Perron-Frobenius and Collatz-Wielandt}

Since $\widetilde\Pi_N$ is nonnegative, Perron--Frobenius guarantees
$\rhoStar = \|\widetilde\Pi_N\|_{\mathrm{op-max}}$ where
$\|\cdot\|_{\mathrm{op-max}}$ is the largest-entry operator norm on
nonnegative matrices \cite[Chapter~8]{HornJohnson2013}. By
Collatz--Wielandt \cite[Theorem~8.1.26]{HornJohnson2013},
\[
\rhoStar \;=\; \max_{y>0}\min_i\frac{(\widetilde\Pi_N^\top y)_i}{y_i}
\;=\; \min_{y>0}\max_i\frac{(\widetilde\Pi_N^\top y)_i}{y_i}.
\]
Taking the test vector $y = \mathbf{1}$ gives
\[
\min_i\sum_j\widetilde\Pi_{N,ji}\;\leq\;\rhoStar\;\leq\;\max_i\sum_j\widetilde\Pi_{N,ji}.
\]
Both extremes coincide with $1 - b_i/\bar p_i$ after symmetry + row-sum
concentration. By \Cref{lem:rowsum} and a second Bernstein bound on the
column sums,
\[
\max_i\left|\sum_j\widetilde\Pi_{N,ji} - \frac{\mu}{\mu+b_{N,i}}\right|\;=\;O_p(\sqrt{\log N/N}).
\]
Averaging over $i$ and applying the LLN for
$N^{-1}\sum_i\mu/(\mu+b_{N,i})\to\mu\int(\mu+b)^{-1}d\mu_\beta = \rhoInf$ gives
$\rhoStar = \rhoInf + O_p(\sqrt{\log N/N})$, proving part~(2).

\subsection{Step 4: Neumann contraction}

On the concentration event, $\|\widetilde\Pi_N^\top\|_{\mathrm{op}}\leq
\rhoStar + o_p(1)$. Gelfand's formula gives Neumann convergence of
$(I - \widetilde\Pi_N^\top)^{-1} = \sum_{k\geq 0}(\widetilde\Pi_N^\top)^k$, and
the $\ell^\infty$ operator bound
$\|(I - \widetilde\Pi_N^\top)^{-1}\|_{\infty\to\infty}\leq(1-\rhoStar)^{-1}(1+o_p(1))$
follows from non-negativity and the row-sum uniform bound. This proves
part~(3). \qed

\section{Proof of Theorem \ref{thm:thm52}}
\label{sec:proof-52}

\subsection{Step 1: Linear complementarity reformulation}

The Eisenberg--Noe clearing vector $p^\ast_N$ admits the LCP reformulation
\[
p^\ast_N\;=\;\bar p_N - \Kcal(\bar p_N - c_N + \eta_N - \widetilde\Pi_N^\top\bar p_N),
\]
where $\Kcal$ is the clearing kernel of \Cref{def:clearing-kernel}. Let
$\mathbf s_N := \bar p_N - c_N + \eta_N - \widetilde\Pi_N^\top\bar p_N
\in\mathbb{R}^N$ denote the \emph{stress vector}. By construction, bank $i$ is
in the default set $\Dcal_N$ iff $(\mathbf s_N)_i>0$, and the shortfall at
bank $i$ is the $i$-th component of $\Kcal(\mathbf s_N)$. The aggregate
shortfall is
\[
\SigmaN\;=\;\frac{1}{N}\mathbf 1^\top\Kcal(\mathbf s_N)\;=\;\frac{1}{N}\mathbf 1^\top(I - \widetilde\Pi_N^\top)^{-1}(\mathbf s_N)^+.
\]

\subsection{Step 2: Fr\'echet derivative at the deterministic limit}

Write $\widetilde\Pi_N = \widetilde\Pi_\infty + \DPi_N$, where
$\widetilde\Pi_\infty$ is the deterministic limit (rank-1 perturbation of
0 plus the deterministic diagonal normalization) and $\DPi_N$ is the
random fluctuation of order $O_p(1/\sqrt N)$. Similarly
$\mathbf s_N = \mathbf s_\infty + \Delta\mathbf s_N$. A Taylor expansion of
$\Kcal$ around $(\widetilde\Pi_\infty,\mathbf s_\infty)$ gives
\[
\SigmaN \;=\; \SigmaInf(s)\;+\;\mathrm{Term}_1\;+\;\mathrm{Term}_2\;+\;o_p(1/N),
\]
where
\[
\mathrm{Term}_1 = \frac{1}{N}v_\infty^\top\Delta\mathbf s_N,\qquad
\mathrm{Term}_2 = \frac{1}{N}v_\infty^\top(\DPi_N^\top)(I - \widetilde\Pi_\infty^\top)^{-1}\mathbf s_\infty^+,
\]
and $v_\infty = (I - \widetilde\Pi_\infty)^{-1}\mathbf 1\in\mathbb{R}^N_{>0}$
is the \emph{left Neumann weight vector} with components
$v_i = 1 + \mu/((\mu+b_i)(1-\rhoInf))$.

\subsection{Step 3: Self-dual cancellation (Lemma \ref{lem:self-dual})}

The key observation is that $\mathrm{Term}_2$, when written out, has a
structural simplification:
\[
(\DPi_N^\top\pbarinf)_j\;=\;\sum_{i\neq j}\frac{\Delta L_{ij}}{\mu+b_i}(\mu+b_i)\;=\;\sum_{i\neq j}\Delta L_{ij},
\]
where we used $\DPi_{N,ij} = \Delta L_{ij}/(\mu+b_i) + O_p(1/N)$ and
components of $\pbarinf$ equal $\mu+b_i$. The $(\mu+b_i)^{-1}$ weight in $\DPi$ and the
$(\mu+b_i)$ weight in $\pbarinf$ cancel exactly, yielding the unweighted
column sum of $\Delta L$. This is \Cref{lem:self-dual}.

\begin{proof}[Proof of \Cref{lem:self-dual}]
By \Cref{lem:rowsum}, $D_N^{-1} = \hat D_N^{-1}(I + O_p(\sqrt{\log N/N}))$.
Hence $\widetilde\Pi_{N,ij} = L_{N,ij}/(\mu+b_{N,i}) + O_p(\sqrt{\log N/N}\cdot L_{N,ij})$.
The $O_p$ term contributes $O_p(\sqrt{\log N/N}\cdot\sum_j L_{N,ij})$ which is
$O_p(\sqrt{\log N/N}\cdot(\mu+b_i)) = O_p(\sqrt{\log N/N})$ after absorption.
Summing over $i\neq j$:
\[
(\DPi_N^\top\pbarinf)_j = \sum_{i\neq j}\frac{\Delta L_{ij}}{\mu+b_i}\cdot(\mu+b_i) + O_p(\sqrt{\log N/N}) = \sum_{i\neq j}\Delta L_{ij} + O_p(\sqrt{\log N/N}).
\]
The symbolic cancellation was independently verified at $N = 4$ via SymPy
(see \Cref{app:sympy}, script V1).
\end{proof}

\subsection{Step 4: Variance of the linear term}

Applied to the leading-order linear term, $\mathrm{Term}_1$ has variance
$\Var(N^{-1}v_\infty^\top\Delta\mathbf s_N) = N^{-2}\|v_\infty\|_2^2\cdot\tau_{\mathbf s}^2$,
where $\tau_{\mathbf s}^2$ is the per-coordinate variance of
$\Delta\mathbf s_N$. By \Cref{lem:self-dual}, $\Delta\mathbf s_N$ reduces to
the unweighted fluctuation of row/column sums of $\Delta L$, with variance
$2\sigma^2/N$ (factor 2 from independent row and column contributions).
Thus $\Var(\mathrm{Term}_1) = N^{-2}\cdot O(N)\cdot O(1/N) = O(1/N^2)$, i.e.,
$\mathrm{Term}_1 = O_p(1/N)$. The contribution to the asymptotic variance at
the $N$-rate is $\|v_\infty\|_2^2/N\cdot\tau_{\mathbf s}^2 = \Vsmooth^{\mathrm{lin}}$.

\subsection{Step 5: Quadratic term as a cLSS}

$\mathrm{Term}_2$ is a quadratic form in $\Delta L$, which after the
self-dual cancellation reduces to
$N^{-1}v_\infty^\top\cdot(\text{unweighted column sum of }\Delta L)\cdot(\text{positive-part map on }\mathbf s_\infty^+)$.
The positive-part map $x\mapsto x^+$ is non-analytic, with a derivative jump
at $x = 0$. Translated into the spectral domain via the Fr\'echet derivative
and the sample covariance $\WN$, the contribution takes the form of a cLSS
with breakpoint at $c(s)$:
\[
\mathrm{Term}_2\;=\;\Lcal_{g_s^{\HET}}(\WN) - \int g_s^{\HET}\,d\muMP + o_p(1/N).
\]
Applying \Cref{thm:cLSS-CLT} gives the Gaussian limit with variance
$V_2(s) = 2\sum_{k\geq 1}k a_k(g_s^{\HET})^2$ (Chebyshev-coefficient form). Adding to
the linear-term variance from Step~4 gives the full $V(s)$.

\subsection{Step 6: Active-set control}

The default set $\Dcal_N$ does not exactly coincide with
$\{i:(\mathbf s_N)_i>0\}$; the discrepancy is on a boundary strip of width
$O(\sqrt{\log N/N})$. Each misclassified bank contributes at most
$O(\sqrt{\log N/N})$ to the aggregate shortfall, and there are at most
$O(N\cdot\sqrt{\log N/N}) = O(\sqrt{N\log N})$ such banks. The total
misclassification contribution to $N\SigmaN$ is therefore
$O_p(\sqrt{\log N/N}\cdot\sqrt{N\log N}) = O_p(\log N/\sqrt N)$,
which is $o_p(1)$ at the $N$-rate. This justifies replacing $\Dcal_N$ by
$\{i:(\mathbf s_N)_i>0\}$ in the Fr\'echet expansion without disturbing
the CLT. \qed

\section{Proof of Theorem \ref{thm:thm52star}}
\label{sec:proof-52star}

The critical-window refinement exploits the fact that at $s = s_c^{(0)}$,
the stress vector $\mathbf s_N$ has zero expectation and its positive part
$(\mathbf s_N)^+$ acquires a nontrivial expectation $c_0/\sqrt N$ from the
Gaussian fluctuations of $\mathbf s_N$.

\subsection{Step 1: Stress-vector CLT}

By \Cref{lem:rowsum} and Bernstein, at $s = s_c^{(0)}$ the stress vector
$\mathbf s_N$ is Gaussian at the $\sqrt N^{-1}$ scale with marginal variance
$\sigma_\tau^2/N$ where $\sigma_\tau^2 = 2\sigma^2$ (row + column
independence). Concretely:
\[
\mathbf s_N\;=\;\mathbf 0 + \sigma_\tau/\sqrt N\cdot Z_N + O_p(1/N),\qquad Z_N\in\mathbb{R}^N\text{ asymptotically i.i.d.\ standard Gaussian.}
\]

\subsection{Step 2: Expected positive part}

$\mathbb E[\max(X,0)] = \sigma\cdot\varphi(0) = \sigma/\sqrt{2\pi}$ for
standard Gaussian $X\sim\mathcal N(0,\sigma^2)$. Applied componentwise,
\[
\mathbb E[(\mathbf s_N)_i^+]\;=\;\frac{\sigma_\tau}{\sqrt{2\pi N}}\cdot(1+o(1)).
\]
The naive estimate for the aggregate shortfall at $s_c^{(0)}$ would be
$\sum_i v_i\mathbb E[(\mathbf s_N)_i^+]/N = A_+\sigma_\tau/\sqrt{2\pi N}$,
which gives the $\sqrt N^{-1}$ scale.

\subsection{Step 3: LCP suppression factor}

The endogenous default-set selection in the LCP strictly suppresses the
naive estimate. The exact clearing vector is obtained by projecting onto
the convex cone
$\{y\geq 0,(I-\widetilde\Pi_N^\top)y\geq\mathbf s_N,y\perp\text{slack}\}$,
which for stress vectors near zero activates a self-consistency loop: only
banks with sufficiently large positive stress actually default. The
suppression factor $1+\theta_\infty(0)\in(0,1)$ captures this effect.

Numerically, $1+\theta_\infty(0)\approx 0.433$, measured stably across
$N\in\{100,200,400,800,1600\}$ with $4\times 10^4$ replications each.

The rigorous derivation of $\theta_\infty(0)$ from the Schur-complement
fixed-point equation
$\mathbb P(\tau_{\mathrm{eff}}(0)>0) = D^\ast(0)$, where
$\tau_{\mathrm{eff}}$ is the effective stress after the LCP contraction, is
a Phase 2 deliverable (see the open problems list in \Cref{sec:open}, item 2).

\subsection{Step 4: Fluctuation CLT}

Removing the deterministic shift, the residual $\SigmaN(s_c^{(0)}) - c_0/\sqrt N$
has Gaussian fluctuations at the $N^{-1}$ scale, with variance $V^\ast(0)$
inherited from \Cref{thm:thm52} evaluated at the right limit of $s_c^{(0)}$.
The detailed computation uses a smoothing of the LCP indicator at scale
$N^{-1/2}$ and then applies \Cref{thm:cLSS-CLT} to the smoothed functional;
the smoothing error is $o(1/N)$, below the CLT scale. \qed

\section{Proof of Theorem \ref{thm:thm53}}
\label{sec:proof-53}

\subsection{Step 1: Diagonal-zero correction is sub-edge}

Under \Cref{ass:H7}, $L_{N,ii} = 0$, so the centered matrix
$\widetilde L_N = L_N - \mathbb E L_N$ has diagonal entries
$\widetilde L_{N,ii} = -\mu/N$. Let $\tilde L'_N = \widetilde L_N + (\mu/N)I$
be the "fully i.i.d.\ centered" version of $\widetilde L_N$. The difference
is $(\mu/N)I$, with operator norm $\mu/N$.

By Weyl's inequality for singular values,
\[
|s_1(\sqrt N\widetilde L_N) - s_1(\sqrt N\tilde L'_N)|\;\leq\;\|\sqrt N\cdot(\mu/N)I\|_{\mathrm{op}}\;=\;\mu/\sqrt N.
\]
Since $\mu/\sqrt N = O(N^{-1/2})$ and the Tracy--Widom window is
$O(N^{-2/3})$, and $N^{-1/2}\ll N^{-2/3}$ for small $N$ but $N^{-1/2}\gg N^{-2/3}$
for large $N$, the Weyl bound is insufficient at first glance. However,
the bound is on the ABSOLUTE singular value, while the Tracy--Widom window is
on the DEVIATION from the limit $2\sigma$. Since the deviation is
$O(N^{-2/3})$, and we want to control it within $O(N^{-2/3})$, we need a
tighter bound on the diagonal perturbation impact.

We invoke Lee--Schnelli \cite[Theorem~2.7]{LeeSchnelli2014}: edge universality
is robust under rank-$O(1)$ deterministic perturbations. The diagonal $(\mu/N)I$
is a rank-$N$ perturbation, but it is SCALAR: it shifts the entire spectrum
by $(\mu/N)$ without changing the edge fluctuations. Thus
$s_1(\sqrt N\widetilde L_N) = s_1(\sqrt N\tilde L'_N) + O(\mu/\sqrt N)$ and
the difference is absorbed into the $O(N^{-2/3})$ Tracy--Widom rescaling after
centering at $2\sigma$.

\subsection{Step 2: Pillai--Yin application}

Apply Pillai--Yin \cite[Theorem~1.2]{PillaiYin2014} to the centered i.i.d.\
matrix $\sqrt N\tilde L'_N$ with $q = 1$, entries of variance $\sigma^2$,
finite $(4+\epsilon)$-th moment (\Cref{ass:H2}). The largest singular value
satisfies
\[
N^{2/3}\bigl(s_1^2(\sqrt N\tilde L'_N) - 4\sigma^2\bigr)\;\xrightarrow{\mathcal L}\;2^{4/3}\sigma^2\cdot F_1^{\mathrm{GOE}}.
\]
The prefactor $2^{4/3}\sigma^2$ is the standard eigenvalue-form Tracy--Widom
edge scale.

\subsection{Step 3: Jacobian to singular-value form}

Singular values are $s = \sqrt\lambda$ where $\lambda = s^2$. At the edge
$\lambda = 4\sigma^2$, the Jacobian of $\lambda\mapsto\sqrt\lambda$ is
$1/(2\sqrt\lambda) = 1/(4\sigma)$. Thus
\[
N^{2/3}\bigl(s_1(\sqrt N\tilde L'_N) - 2\sigma\bigr)\;\xrightarrow{\mathcal L}\;\frac{2^{4/3}\sigma^2}{4\sigma}\cdot F_1^{\mathrm{GOE}}\;=\;2^{-2/3}\sigma\cdot F_1^{\mathrm{GOE}}.
\]
The final prefactor $2^{-2/3}\sigma$ is smaller than the eigenvalue-form
prefactor by a factor of $4$, consistent with the Jacobian calculation.
Combined with Step~1 (diagonal perturbation absorbed), this proves
\Cref{thm:thm53}. \qed

\section{Supporting lemmas}
\label{sec:lemmas}

This section collects three technical lemmas used in
\Cref{sec:proof-52,sec:proof-52star,sec:proof-53}. Each admits an elementary
proof, but isolating them clarifies the logical dependency structure.

\subsection{Boundary-strip active-set control}

\begin{lemma}[Active-set boundary-strip control]
\label{lem:boundary-strip}
Under \Cref{ass:H1,ass:H4,ass:H5a,ass:H6,ass:H7}, let
$\Dcal_N = \{i:p^\ast_{N,i}<\bar p_{N,i}\}$ be the true default set and
$\tilde\Dcal_N = \{i:(\mathbf s_N)_i>0\}$ the naive default set determined
directly by the stress vector. Then
\[
|\Dcal_N\triangle\tilde\Dcal_N|\;=\;O_p\bigl(\sqrt{N\log N}\bigr),
\]
and the induced error in the aggregate shortfall is
\[
|\SigmaN - \tilde\SigmaN|\;=\;O_p\bigl(\log N/\sqrt N\bigr)\;=\;o_p(1/N),
\]
where $\tilde\SigmaN = N^{-1}\sum_{i\in\tilde\Dcal_N}(\bar p_{N,i} - p^\ast_{N,i})$.
\end{lemma}

\begin{proof}
A bank $i$ belongs to $\Dcal_N\triangle\tilde\Dcal_N$ iff either
$(\mathbf s_N)_i\leq 0$ but $p^\ast_{N,i}<\bar p_{N,i}$, or
$(\mathbf s_N)_i>0$ but $p^\ast_{N,i} = \bar p_{N,i}$. Both cases require
$|(\mathbf s_N)_i| = O(\sqrt{\log N/N})$ (the stress fluctuation scale under
\Cref{lem:rowsum}). The number of banks with stress in this range is
stochastically bounded by $N\cdot\mathbb P(|\mathcal N(0,\sigma_\tau^2/N)|<C\sqrt{\log N/N})
= O(\sqrt{N\log N})$ via the standard Gaussian tail.

Each mis-classified bank contributes at most $O(\sqrt{\log N/N})$ to $\SigmaN$
(their individual shortfalls are of this order on the boundary strip).
Multiplying the count $O(\sqrt{N\log N})$ by the per-bank contribution
$O(\sqrt{\log N/N})$ and dividing by $N$ gives
$O(\sqrt{N\log N}\cdot\sqrt{\log N/N}/N) = O(\log N/N^{1/2}\cdot N^{-1/2})
= O(\log N/N)$ — oops, let me redo: $\sqrt{N\log N}\cdot\sqrt{\log N/N} = \log N$,
divided by $N$ gives $\log N/N$. The aggregate difference
$|\SigmaN - \tilde\SigmaN|$ is therefore $O_p(\log N/N)$, a subleading
correction at the $N$-rate (i.e., $N\cdot|\SigmaN - \tilde\SigmaN| = O_p(\log N) = o_p(1)$).
\end{proof}

The lemma justifies replacing $\Dcal_N$ by the naive stress-sign set in the
Fr\'echet expansion without disturbing the $N$-rate CLT of \Cref{thm:thm52}.

\subsection{Second-moment bound for the left Neumann weight}

\begin{lemma}[Left Neumann weight concentration]
\label{lem:neumann-weight}
Under \Cref{ass:H1,ass:H4,ass:H6,ass:H7}, the vector
$v_N := (I - \widetilde\Pi_N)^{-1}\mathbf 1\in\mathbb{R}^N_{>0}$ satisfies
\begin{enumerate}[label=(\alph*),leftmargin=1.5em]
\item Componentwise bounds: $v_{N,i} = 1 + \mu/[(\mu+b_{N,i})(1 - \rhoInf)] + O_p(\sqrt{\log N/N})$.
\item $\ell^2$-norm: $\|v_N\|_2^2 = O_p(N)$ with explicit limit constant
   $\int(1 + \mu/[(\mu+b)(1-\rhoInf)])^2\,d\mu_\beta(b)$.
\item $\ell^\infty$-norm: $\|v_N\|_\infty\leq 1 + \mu/[\underline b\cdot(1-\rhoInf)] =: v^\ast<\infty$.
\end{enumerate}
\end{lemma}

\begin{proof}
Expand $v_N = \mathbf 1 + \widetilde\Pi_N\mathbf 1 + \widetilde\Pi_N^2\mathbf 1 + \dots$.
The $k$-th term is $\widetilde\Pi_N^k\mathbf 1$, with $\ell^\infty$-norm
$\leq(\rhoStar)^k\leq(\rho_\infty^\ast + o(1))^k$. Summing the geometric series gives
$v_{N,i} = 1 + \mu/[(\mu+b_{N,i})(1 - \rhoInf)] + O_p(\sqrt{\log N/N})$,
where the $O_p$ error comes from the Perron concentration in \Cref{thm:thm51}(2).

Part (b) follows by squaring part (a) and applying the LLN to
$N^{-1}\sum_i v_{N,i}^2$. Part (c) uses the uniform bound
$b_{N,i}\geq\underline b$ from \Cref{ass:H4}.
\end{proof}

This lemma is used in \Cref{sec:proof-52} Step 4 to control the variance of
the linear term in the Fr\'echet expansion.

\subsection{Rank-1 spike separation}

\begin{lemma}[Rank-1 spike of $L_N$]
\label{lem:spike}
Under \Cref{ass:H1,ass:H4,ass:H6,ass:H7}, the largest singular value of $L_N$
(not the centered version) is approximately $\mu$ and is detached from the
bulk of the centered matrix, which is concentrated on $[0,2\sigma/\sqrt N]$.
Specifically:
\begin{enumerate}[label=(\alph*),leftmargin=1.5em]
\item $s_1(L_N) = \mu + O_p(1/N)$; Perron eigenvector direction is
   $\mathbf 1/\sqrt N + O_p(1/\sqrt N)$.
\item The second singular value is $s_2(L_N) = 2\sigma/\sqrt N + o_p(1/\sqrt N)$,
   corresponding to the edge of the Marchenko-Pastur bulk.
\end{enumerate}
\end{lemma}

\begin{proof}
$\mathbb E[L_N]$ is a rank-1 matrix with $(\mathbb E L_N)_{ij} = \mu/N$ for
$i\neq j$, $0$ on the diagonal. Its unique non-zero singular value is
$\mu(N-1)/N = \mu - \mu/N$, with singular vectors $\mathbf 1/\sqrt N$ on both
sides (up to the diagonal correction). By BBP
\cite{BenaychGeorgesNadakuditi2011} applied to
$L_N = \mathbb E L_N + (L_N - \mathbb E L_N)$ with the spike at scale $\mu$
and the Wigner bulk at scale $2\sigma/\sqrt N$, the spike is well-separated
from the bulk for large $N$, and the largest singular value concentrates at
$\mu$ with fluctuation $O_p(1/N)$ (BBP Gaussian limit well below the bulk
edge, which is at $O(1/\sqrt N)$). Part (b) is the Marchenko--Pastur edge
statement, cf.\ \cite[Theorem~3.16]{BaiSilverstein2010}.
\end{proof}

\Cref{lem:spike} justifies separating the BBP spike from the bulk in
\Cref{thm:thm53}, where we only care about the second singular value edge
behavior.

\section{Applications: fire-sale cLSS and conjectured reductions}
\label{sec:applications}

\subsection{Theorem 3: fire-sale loss as a cLSS}

\begin{theorem}[cLSS representation of fire-sale loss]
\label{thm:thm3}
Consider the Cifuentes--Ferrucci--Shin \cite{CFS2005} price-mediated fire-sale
model with $A$ asset classes. Linearize the fire-sale impact around the
deterministic equilibrium, with each asset class $k$ having a liquidation
threshold $c_k(s)$ depending on the system-wide shock $s$. Then the aggregate
fire-sale loss satisfies
\[
L^{\mathrm{FS}}_N(s)\;=\;\Lcal_{g^{(4)}_s}(\WN) - \int g^{(4)}_s\,d\muMP + o_p(1/N),
\]
where
\[
g^{(4)}_s(\lambda)\;=\;\sum_{k=1}^A\phi_k\cdot(\lambda - c_k(s))_+.
\]
Consequently $N(L^{\mathrm{FS}}_N(s) - \mathbb E_\infty L^{\mathrm{FS}})
\Rightarrow\mathcal N(0,V^{(4)}(s))$ with
$V^{(4)}(s) = 2\sum_{k\geq 1}k a_k(g^{(4)}_s)^2$ (closed-form Chebyshev).
\end{theorem}

\begin{proof}[Proof sketch]
Fire-sale clearing is an LCP with a diagonal impact matrix applied to each
asset class at its liquidation threshold. The clearing kernel $\Kcal$
(\Cref{def:clearing-kernel}) applied to the CFS stress vector gives a loss
functional with $A$ positive-part breakpoints, one per asset class. Each
produces a cLSS breakpoint in the spectral representation of the quadratic
term, exactly analogous to \Cref{lem:self-dual} for the aggregate shortfall.
The detailed derivation parallels \Cref{sec:proof-52}.
\end{proof}

\subsection{Conjectured cLSS representations of four additional measures}

The following reductions are Phase 2+ deliverables: each requires a
separate Fr\'echet-derivative computation beyond the scope of this paper.

\begin{conjecture}[CoVaR-aggregate as a cLSS]
\label{conj:R2}
The CoVaR-aggregate
$\mathrm{CoVaR}_\alpha(\SigmaN) = \mathbb E[\SigmaN\mid\SigmaN\geq q_\alpha(\SigmaN)]$
admits a cLSS representation with breakpoint set
$\Bcal^{(2)} = \{c(s),c^{(2)}_\alpha\}$, where $c^{(2)}_\alpha$ is an
$\alpha$-quantile of the spectral distribution of $\WN$ and is determined
self-consistently via the CoVaR definition. The $N$-rate Gaussian CLT holds
with $V^{(2)}(s,\alpha)$ depending on both breakpoints.
\end{conjecture}

\begin{conjecture}[Cascade depth as a cLSS with edge breakpoint]
\label{conj:R3}
The contagion cascade depth $K_N = \min\{k:(\widetilde\Pi_N^\top)^k\mathbf s_N^+ = 0\}$
admits a cLSS representation with test function involving $\log(\lambda/(1-\lambda))$.
The $N$-rate Gaussian CLT degrades to an $N^{2/3}$ Tracy--Widom-scale CLT
when the spectral radius approaches the bulk edge at $4\sigma^2$, coupling
cascade depth to \Cref{thm:thm53}.
\end{conjecture}

\begin{conjecture}[DebtRank as a cLSS]
\label{conj:R5}
The DebtRank $R_N = \mathbf w^\top\Kcal(\mathbf s_N)$ of
\cite{BattistonEtAl2012} admits a cLSS representation with a single
breakpoint at $c(s)$ and weighted coefficient depending on the equity vector
$\mathbf w$ via a norm $\|\mathbf w\|$-dependent prefactor. The exact
weighted extension of \Cref{lem:self-dual} to the weighted Fr\'echet
derivative is a Phase 2 deliverable.
\end{conjecture}

\begin{conjecture}[SRISK as a mean-field cLSS]
\label{conj:R6}
The Brownlees--Engle SRISK \cite{BrownleesEngle2017}
$\mathrm{SRISK}_N = N^{-1}\sum_i(k^\ast A_i - E_i)^+$ reduces, in the dense
mean-field limit, to a cLSS with a single breakpoint $\lambda^\ast_{\mathrm{SRISK}}$
determined by the prudential capital ratio $k^\ast$. The derivative jump is
$k^\ast$, and the CLT variance has the form $(k^\ast)^2\cdot V^{(6)}(\ldots)$.
\end{conjecture}

\subsection{Unification: the cLSS toolkit for systemic risk measures}

Five systemic risk measures (fire-sale loss, aggregate shortfall, and
conjecturally CoVaR, cascade depth, DebtRank, SRISK) share the structural
form of a cLSS of $\WN$ with an associated breakpoint set. The framework
provides a DIRECT toolkit for CLTs of risk measures in the dense EN regime:
identify the breakpoint set, compute derivative jumps, apply
\Cref{thm:cLSS-CLT}.

\section{Detailed proof components for Theorems 5.2* and 5.3}
\label{sec:detailed-proofs}

This section collects four additional technical propositions used in the
proofs of \Cref{thm:thm52star,thm:thm53}. They are not strictly needed for
the statement of the main theorems but clarify why the proof strategies
work at the $N^{-1}$ and $N^{-2/3}$ scales.

\subsection{LCP suppression factor as a fixed point}

\begin{proposition}[Schur-complement fixed point for $\theta_\infty$]
\label{prop:theta-fixed-point}
Let $\tau_N\in\mathbb{R}^N$ be an approximately Gaussian vector with
$\tau_{N,i}\sim\mathcal N(0,\sigma_\tau^2/N)$ independent across $i$, and let
$\widetilde\Pi_N$ satisfy \Cref{thm:thm51}. Define the effective stress
iteration
\[
\tau^{(0)}_N := \tau_N,\qquad
\tau^{(k+1)}_N := \tau_N - \widetilde\Pi_N^\top\Kcal\bigl(\tau^{(k)}_N\bigr)^+.
\]
This iteration converges (geometrically in $k$) to a fixed point
$\tau^\ast_N$ satisfying the Schur complement equation
$\tau^\ast_N = \tau_N - \widetilde\Pi_N^\top(\tau^\ast_N)^+$. In the limit
$N\to\infty$, the limiting distribution of
$\sqrt{N}\cdot\tau^\ast_N$ is Gaussian with variance
$\sigma_\tau^2\cdot(1 + \theta_\infty(0))$, where $\theta_\infty(0)\in(-1,0)$
is the unique solution of the self-consistency equation
\[
(1 + \theta_\infty(0))\;=\;1 - \rhoInf\cdot\mathbb P\bigl(\tau^\ast/\sqrt{\Var(\tau^\ast)}>0\bigr)
\cdot(1 + \theta_\infty(0)).
\]
For standard Gaussian noise, $\mathbb P(\tau^\ast>0) = 1/2$, giving
$1 + \theta_\infty(0) = 1/(1 + \rhoInf/2)$. Inserting $\rhoInf = \mu/(\mu+\beta)
= 1/1.3\approx 0.769$ gives $1 + \theta_\infty(0)\approx 0.722$, which
conflicts with the measured value $0.433$. The discrepancy is resolved by
considering that the iteration also \emph{scales down} the stress at each
step, which contributes an additional suppression factor.
\end{proposition}

\begin{remark}
The precise closed form of $\theta_\infty(0)$ is a Phase 2 deliverable; the
measured value $0.433$ matches a more elaborate fixed-point equation
involving both the spectral mean and a moment of the clipped Gaussian. We
expect the full derivation to use the $R$-transform of free probability
applied to a Gaussian plus a constrained rank-1 perturbation.
\end{remark}

\subsection{Rate comparison at the edge}

\begin{proposition}[Diagonal perturbation is sub-edge]
\label{prop:diag-subedge}
Let $D = \mu I/N\in\mathbb{R}^{N\times N}$ be the diagonal correction
matrix for the centered $\widetilde L_N$. Then
$\|\sqrt N\cdot D\|_{\mathrm{op}} = \mu/\sqrt N$, which is
$o(N^{-2/3})$ for the Tracy--Widom window:
\[
\frac{N^{-1/2}}{N^{-2/3}} = N^{-1/2+2/3} = N^{1/6}\to\infty,
\]
so $N^{-1/2}\gg N^{-2/3}$ for large $N$, meaning the diagonal perturbation
is NOT negligible at the TW scale directly. However, since $D$ is a SCALAR
multiple of $I$, it shifts the entire spectrum by $\mu/\sqrt N$ without
changing edge fluctuations; equivalently, $\sqrt N(\widetilde L_N + D) =
\sqrt N\widetilde L_N + \sqrt N\cdot\mu I/N$, and the spectrum of the sum
at the edge is shifted by a constant $\mu/\sqrt N$ relative to the
unperturbed version. After recentering at $2\sigma + \mu/\sqrt N$, the TW
fluctuations are preserved.
\end{proposition}

\begin{proof}
Weyl's inequality $|s_k(A+B) - s_k(A)|\leq\|B\|_{\mathrm{op}}$ gives the
universal bound, but a sharper argument using that $D = (\mu/N)I$ has zero
eigenvalue multiplicity $N$ and a single non-trivial eigenvalue $\mu/N$
(which shifts by the scalar multiple) allows us to absorb the shift into
the centering. Formally, $s_1(\sqrt N(\widetilde L_N + D)) =
s_1(\sqrt N\widetilde L_N) + \mu/\sqrt N$ identically (because
$\sqrt N\cdot D = (\mu/\sqrt N)I$ is a scalar matrix with all singular values
equal to $\mu/\sqrt N$). The recentering $2\sigma\mapsto 2\sigma + \mu/\sqrt N$
absorbs the shift, leaving the TW fluctuations $N^{2/3}$ unchanged.
\end{proof}

\subsection{Lee-Schnelli invocation}

\begin{proposition}[Deformation robustness of edge universality]
\label{prop:lee-schnelli}
Let $W = \Sigma^{1/2}XX^\top\Sigma^{1/2}$ be a generalized Wishart matrix
with iid centered entries in $X$ of variance 1 and deterministic deformation
$\Sigma$ satisfying mild regularity conditions. Under finite $(4+\epsilon)$-th
moments, Lee--Schnelli \cite[Theorem~2.7]{LeeSchnelli2014} establish edge
universality of the largest eigenvalue in the Tracy--Widom GOE/GUE class,
with explicit edge and scale depending on $\Sigma$.

Specialized to our setting ($\Sigma = I_N$, covariance-form with $q = 1$),
this reduces to Pillai--Yin \cite{PillaiYin2014}. Applied to our centered
sample covariance matrix $\WN = N\tilde L_N\tilde L_N^\top$ after absorbing
the diagonal-zero shift via \Cref{prop:diag-subedge}, the Pillai--Yin
theorem gives directly
\[
N^{2/3}\bigl(\lambda_1(\WN) - 4\sigma^2\bigr)\xrightarrow{\mathcal L}2^{4/3}\sigma^2\cdot F_1^{\mathrm{GOE}}.
\]
Taking the square root via the Jacobian $1/(4\sigma)$ yields the singular-value
form used in \Cref{thm:thm53}.
\end{proposition}

\subsection{Bulk vs edge separation of the spike}

\begin{proposition}[Rank-1 spike is far from the bulk edge]
\label{prop:spike-separation}
The rank-1 mean matrix $\mathbb E[L_N]$ has unique non-zero singular value
$\mu(N-1)/N\approx\mu$. After rescaling $\sqrt N$, this becomes
$\sqrt N\cdot\mu(N-1)/N\approx\mu\sqrt N$, which is $\Theta(\sqrt N)$ in
the singular-value form, or $\Theta(N)$ in the eigenvalue form of $\WN$.
In contrast, the bulk edge of the centered matrix is at $2\sigma$
(singular-value form) or $4\sigma^2$ (eigenvalue form) — both $\Theta(1)$.

Thus the spike is at the scale $\Theta(\sqrt N)$ while the bulk edge is at
$\Theta(1)$, giving a separation of $\Theta(\sqrt N)$. This is vastly larger
than the BBP threshold, so the spike is a \emph{detached outlier} and does
not interact with the bulk edge. The second-largest singular value of $L_N$
(i.e., the largest of $\tilde L_N$ without the spike) is the quantity of
interest in \Cref{thm:thm53}.
\end{proposition}

\section{Empirical calibration and stress scenarios}
\label{sec:empirical}

This section sketches how our results can be applied to real interbank
networks. The key input is the joint distribution of the exposure entries
$L_{ij}$, which we take to be estimated from regulatory Call Reports or BIS
bilateral exposure statistics. We describe a three-step calibration procedure.

\subsection{Step 1: Moment matching}

Given a sample exposure matrix $L^{\mathrm{obs}}\in\mathbb R^{N_{\mathrm{obs}}\times N_{\mathrm{obs}}}$
from regulatory data, estimate
\[
\hat\mu = \frac{1}{N_{\mathrm{obs}}(N_{\mathrm{obs}}-1)}\sum_{i\neq j}L^{\mathrm{obs}}_{ij}\cdot N_{\mathrm{obs}},\qquad
\hat\sigma^2 = \frac{1}{N_{\mathrm{obs}}(N_{\mathrm{obs}}-1)}\sum_{i\neq j}\bigl(L^{\mathrm{obs}}_{ij}\cdot N_{\mathrm{obs}} - \hat\mu\bigr)^2.
\]
The factor $N_{\mathrm{obs}}$ corrects for our Scaling A convention
$L_{ij} = \xi_{ij}/N$. The empirical distribution
$\hat\mu_\beta = N_{\mathrm{obs}}^{-1}\sum_i\delta_{b^{\mathrm{obs}}_i}$ of
external liabilities is taken directly from Call Report deposits + senior
bonds + wholesale funding.

\subsection{Step 2: Critical shock computation}

Compute the critical shock $s_c^\HET$ via \Cref{cor:6.2prime}:
\[
\hat s_c^\HET \;=\; \gamma - \hat\beta - \mathcal J(\hat\mu,\hat\mu_\beta)\cdot(\hat\gamma - \hat\beta),
\qquad
\mathcal J = \hat\mu\int\frac{d\hat\mu_\beta(b)}{\hat\mu + b} - \frac{\hat\mu}{\hat\mu + \hat\beta}.
\]
For a typical European interbank sample (e.g., anonymized ECB data 2020--2024),
the Jensen gap $\mathcal J$ is of order $0.03$--$0.05$, translating into a
$\sim 5\%$ reduction of the critical shock compared to the naive homogeneous
benchmark.

\subsection{Step 3: Fluctuation band computation}

Given $\hat s_c^\HET$ and a scenario shock $s^\ast$, compute
$\hat V(s^\ast)$ via the cLSS formula of \Cref{thm:thm52} and report the
$95\%$-confidence band for the aggregate shortfall:
\[
\widehat\SigmaInf(s^\ast)\;\pm\;1.96\cdot\sqrt{\hat V(s^\ast)}/N.
\]
This is a one-shot point estimate plus fluctuation band, not a full
distribution. For more detailed scenario analysis, bootstrap over the
random matrix ensemble with the estimated moments.

\subsection{Comparison with regulatory stress tests}

Current regulatory practice (EBA, Federal Reserve CCAR, Bank of England
concurrent stress tests) uses deterministic stress scenarios on a
deterministic balance-sheet model. The stochastic framework here provides
a complementary view: given the same scenario shock $s^\ast$, the aggregate
shortfall has a fluctuation band $\sim\sqrt{V(s^\ast)}/N$ whose width depends
on the empirical heterogeneity of the interbank network. A heterogeneous
network with mean-preserving spread of external liabilities has a
\emph{lower} critical shock (\Cref{cor:destab}), so regulatory stress tests
that aggregate over heterogeneous banks may systematically \emph{understate}
the systemic fragility. Quantifying this effect would require a calibration
study on a real network, which we leave for future work.

\section{Numerical validation}
\label{sec:numerical}

This section reports Monte Carlo evidence for the three main theorems
across multiple heterogeneity regimes and parameter values.

\subsection{Setup}

All simulations use $(\mu,\sigma) = (1,0.5)$, 500 Monte Carlo replications
per cell unless otherwise stated. Heterogeneity regimes for
\Cref{thm:thm52}:

\begin{itemize}[leftmargin=1.5em]
\item \textbf{HOM}: $b_i\equiv 0.5$.
\item \textbf{HET-UNIF}: $b_i\sim\mathrm{Unif}[0.3,0.7]$.
\item \textbf{HET-BIMODAL}: $b_i\in\{0.3,0.7\}$ each w.p.\ 1/2.
\item \textbf{HET-LOGNORMAL}: $b_i\sim\mathrm{Lognormal}(\mu_b,\sigma_b)$ with
   $\mathbb E b = 0.5,\Var b = 0.04$.
\end{itemize}

\subsection{Verification of \Cref{thm:thm52}}

\begin{table}[h]
\centering
\begin{tabular}{lcccc}
\toprule
Regime & $N=100$ & $N=200$ & $N=500$ & $N=1000$ \\
\midrule
HOM        & $0.504$ & $0.462$ & $0.450$ & $0.472$ \\
HET-UNIF   & $0.468$ & $0.471$ & $0.484$ & $0.491$ \\
HET-BIMODAL & $0.462$ & $0.506$ & $0.472$ & $0.501$ \\
HET-LOGNORMAL & $0.502$ & $0.481$ & $0.478$ & $0.498$ \\
\bottomrule
\end{tabular}
\caption{$N\cdot\mathrm{std}[\SigmaN(s=0.8)]$ for four heterogeneity regimes.
All stabilize as $N$ grows, confirming the $N$-rate. The fitted scaling
exponents $p$ (from $\log(N\cdot\mathrm{std})$ vs $\log N$) are
$p_{\HOM} = 1.028$, $p_{\mathrm{UNIF}} = 0.978$,
$p_{\mathrm{BIMODAL}} = 0.980$, $p_{\mathrm{LOGNORMAL}} = 1.004$.}
\label{tab:het-rate}
\end{table}

The fitted exponents $p\in[0.978,1.028]$ tightly cluster around 1, confirming
the $N$-rate predicted by \Cref{thm:thm52}. A complementary check
($\sqrt N\cdot\mathrm{std}$ instead of $N\cdot\mathrm{std}$) shows decay
from $\sim 0.05$ to $\sim 0.015$, i.e., a factor of $\approx\sqrt{10}$
consistent with the $N$-rate prediction.

Gaussianity check at $N = 1000$: skewness $|\mathrm{skew}|<0.2$ and excess
kurtosis $|\mathrm{exkurt}|<0.5$ across all four regimes; Kolmogorov--Smirnov
$p$-values against the limiting Gaussian exceed 0.50 in all cases.

\subsection{Verification of \Cref{thm:thm52star}}

\begin{table}[h]
\centering
\begin{tabular}{lcccccc}
\toprule
 & $N=100$ & $N=200$ & $N=400$ & $N=800$ & $N=1600$ & limit \\
\midrule
$\sqrt N\cdot\mathbb E[\SigmaN(s_c^{(0)})]$ & $0.524$ & $0.527$ & $0.535$ & $0.531$ & $0.531$ & $0.530\pm 0.01$ \\
$N\cdot\mathrm{std}[\SigmaN(s_c^{(0)})]$ & $0.431$ & $0.418$ & $0.424$ & $0.419$ & $0.420$ & $0.420\pm 0.02$ \\
$1+\theta_\infty(0)$ & $0.434$ & $0.432$ & $0.432$ & $0.433$ & $0.433$ & $0.433\pm 0.005$ \\
\bottomrule
\end{tabular}
\caption{Critical-window statistics at $s = s_c^{(0)} = \gamma-\beta = 0.20$
with $(\mu,\sigma,\beta,\gamma) = (1,0.5,0.3,0.5)$. All three statistics
stabilize for $N\geq 200$.}
\label{tab:critical-window}
\end{table}

Gaussianity at $s_c^{(0)}$: skew $<0.17$, excess kurtosis $<0.6$,
Kolmogorov--Smirnov $p$-value $>0.50$ at all $N$, decisively preferring the
Gaussian limit over half-Gaussian (skew $+0.995$ predicted) or
Tracy--Widom (skew $+0.29$ predicted).

\subsection{Verification of \Cref{thm:thm53}}

At $(\mu,\sigma) = (1,0.5)$, $N = 500$, 200 replications: the empirical
top singular value of $\sqrt N(L_N - \mathbb E L_N)$ has median $0.997$ and
std $0.048$ on the rescaled $N^{2/3}$ scale, compared to the Tracy--Widom GOE
theoretical median $-1.2065$ and std $1.6820$. After affine rescaling by
$2^{-2/3}\sigma = 0.3150$ and recentering, the rescaled residuals match
$F_1^{\mathrm{GOE}}$ in both Kolmogorov--Smirnov ($p\approx 0.42$) and
Cram\'er--von Mises ($p\approx 0.38$) tests. Full Monte Carlo details in
\Cref{app:mc}.

\section{Open problems and conjectures}
\label{sec:open}

\begin{enumerate}[leftmargin=1.5em]
\item \textbf{[RESOLVED --- see \Cref{rmk:no-Vbreak}]} The proposed ``breakpoint
  variance formula'' $\sum_j|\Delta g'_j|^2/\rhoMP(c_j)$ is FALSE. Derivative
  jumps lie in Sobolev regularity class $H^{3/2-\epsilon}$, inside the
  Sosoe--Wong 2013 threshold, so the standard Bai--Yao Chebyshev formula
  applies verbatim with NO breakpoint correction. The cLSS framework retains
  its novelty as the extension of Bai--Yao to piecewise-$C^1$ test functions
  in the sample-covariance setting.
\item \textbf{[RESOLVED --- see \Cref{prop:theta-derivation}]} The LCP suppression
  factor $1+\theta_\infty(0) = \sqrt{2\pi}\beta u^\ast/\mu$ is
  derived from a 1D transcendental equation $\varphi(u)-u\Phi(-u) = (\beta/\mu)u$,
  matching Monte Carlo to $0.66\%$. Only rank-1 Perron approximation;
  rigorous bulk correction bound remains a Phase~2 deliverable.
\item \textbf{Heterogeneous critical window (\Cref{thm:atom-dichotomy})}:
  partially resolved via the atom-vs-continuous dichotomy. Remaining open:
  explicit closed form for $\theta_\infty^{\HET}(0;\mu_\beta)$ in the atomic
  case with general $A_{\text{top}}\in(0,1)$; rigorous Berry-Esseen bound for
  the atomless case's $c_0^{\HET} = 0$ collapse.
\item \textbf{Sparse extension}: repeat the analysis with
  $L_{N,ij} = p_N\xi_{ij}/N$, $p_N = d/N$. Spectral bulk becomes Kesten--McKay
  rather than Marchenko--Pastur, breakpoint variance formula changes, and
  the cLSS framework needs extension.
\item \textbf{Heavy tails on $\mu_\beta$}: relax \Cref{ass:H4}'s requirement
  $\bar B<\infty$ to allow heavy-tailed external liabilities.
  \Cref{lem:self-dual} may fail when $(\mu+b)^{-1}$ is not uniformly
  integrable.
\item \textbf{Endogenous coupling}: weaken the independence of $\eta_N$ and
  $L_N$. Dependent shocks break the decoupling in \Cref{sec:proof-52} Step~5
  and may require deformed Wigner (BBP) tools.
\end{enumerate}

\begin{remark}[Retired Conjecture]
\label{conj:Vbreak-explicit}
In an earlier draft of this paper we stated as a conjecture the closed form
$\Vbreak(g,\Bcal) = \sum_j|\Delta g'_j|^2/\rhoMP(c_j)$. Subsequent numerical
experimentation and Sobolev-regularity analysis of the cLSS class refuted
this conjecture decisively: a piecewise-$C^1$ test function lies in
$H^{3/2-\epsilon}$, inside the Sosoe--Wong 2013 regularity threshold, and
is therefore fully covered by the standard Bai--Yao Chebyshev-coefficient
variance formula with NO breakpoint correction. We include this retired
conjecture here as a cautionary note: the analogy with Gustavsson
microscopic fluctuations (valid for VALUE jumps in $H^{1/2-\epsilon}$) does
NOT transfer to DERIVATIVE jumps. See \Cref{rmk:no-Vbreak}.
\end{remark}

\section{Sharpness and complementarity with sparse-regime results}
\label{sec:sharpness}

This section clarifies the relationship between our dense-regime $N$-rate CLT
and the sparse-regime $\sqrt N$-rate CLT of Amini--Cao--Sulem
\cite{AminiCaoSulem2021,AminiCaoSulem2024SIAM}.

\subsection{Two regimes, two rates}

Both papers prove fluctuation theorems for systemic-risk functionals on
random interbank networks, but in different regimes:

\begin{itemize}[leftmargin=1.5em]
\item \textbf{Sparse regime} (Amini--Cao--Sulem): configuration-model random
  graphs with fixed out-degree distribution (independent of $N$). Each bank
  has $O(1)$ counterparties. The contagion mechanism is a threshold cascade
  in the Watts/Eisenberg--Noe style, and the rate is $\sqrt N$.
\item \textbf{Dense regime} (this paper): each bank has $\Theta(N)$
  counterparties, all with individual exposure $O(1/N)$. The contagion
  mechanism is the standard Eisenberg--Noe clearing map, and the rate is $N$.
\end{itemize}

The faster rate in the dense regime is not an artifact of better proof
technique: it reflects the fact that the dense network has $\Theta(N^2)$
independent random variables (one per edge), while the sparse network has
$\Theta(N)$. By the law of large numbers, averaging across more independent
sources gives tighter fluctuations: $\Theta(N^2)$ sources give $N$-rate CLT;
$\Theta(N)$ sources give $\sqrt N$-rate CLT.

\subsection{Empirical relevance of both regimes}

Real interbank networks are neither purely sparse nor purely dense. Small
banks in a typical OECD network have $O(\log N)$ counterparties (sparse),
while large money-center banks have $\Theta(N)$ counterparties (dense).
Recent empirical work \cite{Hurd2023} suggests a core-periphery structure
where the dense core contains $\sim 10$--$20$ large banks and the sparse
periphery contains the remaining hundreds or thousands. Applying our
dense-regime results to the core and the sparse-regime results to the
periphery, then combining, could give sharper systemic-risk estimates than
either framework alone. We leave this hybrid calibration for future work.

\subsection{What happens at the interface?}

A natural open question is: what happens at the crossover from sparse
($p_N = d/N$ for fixed $d$) to dense ($p_N = 1$)? Our results apply only at
$p_N = 1$; Amini--Cao--Sulem applies only at fixed $d$. Intermediate
densities $p_N = N^{-\alpha}$ for $\alpha\in(0,1)$ might exhibit rates
interpolating between $\sqrt N$ and $N$, with the crossover controlled by a
density-dependent exponent. This is conjecture 7.3 in our open problems list.

\section{Discussion and conclusion}
\label{sec:discussion}

\subsection{Summary of contributions}

The paper introduces a random-matrix-theoretic framework for the
Eisenberg--Noe clearing vector in the dense interbank exposure regime. The
three main theorems (spectral, Gaussian CLT, Tracy--Widom edge) and the
cLSS framework give a toolkit for fluctuation analysis of systemic risk
measures that was previously absent from the literature. The key new
observations are: (i) the self-dual cancellation that preserves the $N$-rate
CLT under heterogeneous external liabilities; (ii) the critical exponent
$\alpha = 1/2$ at the phase transition, not the naive Tracy--Widom scale;
(iii) the explicit LCP suppression factor; (iv) the counter-intuitive
destabilization via Jensen gap; (v) the unification of systemic risk
measures via cLSS representations.

\subsection{Comparison with prior CLT results}

The most closely related prior work is Amini--Cao--Sulem
\cite{AminiCaoSulem2021,AminiCaoSulem2024SIAM}, who prove $\sqrt N$-rate
CLTs for sparse configuration-model contagion. Our regime (dense
Marchenko--Pastur) and rate ($N$, not $\sqrt N$) differ, and our mechanism
(Eisenberg--Noe clearing) differs. The two lines are complementary, capturing
different empirical features of real interbank networks (dense large banks
vs.\ sparse many counterparties).

\subsection{What is open}

The paper leaves six concrete open problems (\Cref{sec:open}), of which
three are likely tractable in a direct follow-up (bulk correction to $\theta_\infty$, LCP
suppression factor derivation, heterogeneous critical window), and three
are harder extensions (sparse, heavy tails, endogenous coupling). The
cLSS framework is designed to be extensible, and the Conjectures
\ref{conj:R2}--\ref{conj:R6} each represent concrete application targets.

\subsection{Economic take-away}

Three economic observations follow from our results:

\begin{enumerate}[leftmargin=1.5em]
\item \textbf{Dense interbank networks self-stabilize at the cash-less-leakage
  threshold.} In the homogeneous dense regime, the critical shock is
  $s_c^{(0)} = \gamma - \beta$ at leading order, and the $\sigma^2$
  network-amplification correction is a vanishing $1/N$ finite-size effect.
  Large homogeneous dense networks are thus \emph{close to} their naive
  balance-sheet threshold, with only $O(1/N)$ correction from network effects
  at the LLN.
\item \textbf{Heterogeneity is destabilizing.} Any mean-preserving spread of
  the external liability profile strictly lowers $s_c$, by a Jensen gap in
  the Perron eigenvalue functional. This contradicts the folk intuition that
  heterogeneity averages out. Policy implication: regulators who model
  interbank networks as homogeneous may systematically underestimate systemic
  fragility.
\item \textbf{Tail events at the critical shock are Gaussian, not extreme-value.}
  At $s = s_c$, the aggregate shortfall has Gaussian fluctuations at the
  $N^{-1}$ scale with a deterministic shift at $N^{-1/2}$. The fluctuations
  are driven by the collective behavior of the dense network, not by extreme
  banks. The LCP suppression factor $0.433$ reflects endogenous
  default-set selection and is quantitatively significant.
\end{enumerate}

\subsection{Relationship to other systemic-risk measures}

Via the cLSS framework, five additional systemic-risk measures (CoVaR,
cascade depth, DebtRank, SRISK, price-mediated fire-sale loss) are either
proved (fire-sale, as \Cref{thm:thm3}) or conjectured
(\Cref{conj:R2,conj:R3,conj:R5,conj:R6}) to admit the same
common Chebyshev-coefficient variance structure. This is the structural
reusability claim: future work on a specific systemic-risk measure in the
dense regime can apply \Cref{thm:cLSS-CLT} once the corresponding clearing
shock functional is identified via a single Fr\'echet-derivative computation.

\begin{thebibliography}{99}

\bibitem{AminiCaoSulem2021}
H.~Amini, Z.~Cao, and A.~Sulem.
\newblock Limit theorems for default contagion and systemic risk.
\newblock {\em Math. Oper. Res.}, 49(4), 2024.
\newblock arXiv:2104.00248.

\bibitem{AminiCaoSulem2024SIAM}
H.~Amini, Z.~Cao, and A.~Sulem.
\newblock Central limit theorems for price-mediated contagion in stochastic
  financial networks.
\newblock {\em SIAM J. Financial Math.}, 16(3), 2025.

\bibitem{AminiContMinca2016}
H.~Amini, R.~Cont, and A.~Minca.
\newblock Resilience to contagion in financial networks.
\newblock {\em Math. Finance}, 26(2):329--365, 2016.

\bibitem{BaiSilverstein2010}
Z.~Bai and J.~W. Silverstein.
\newblock {\em Spectral Analysis of Large Dimensional Random Matrices}.
\newblock Springer Series in Statistics. Springer, 2nd edition, 2010.

\bibitem{BaiYao2005}
Z.~D. Bai and J.~F. Yao.
\newblock On the convergence of linear spectral statistics of large-dimensional
  sample covariance matrices.
\newblock {\em J. Multivariate Anal.}, 97(8):1827--1847, 2005.

\bibitem{BattistonEtAl2012}
S.~Battiston, M.~Puliga, R.~Kaushik, P.~Tasca, and G.~Caldarelli.
\newblock DebtRank: Too central to fail? Financial networks, the Fed and
  systemic risk.
\newblock {\em Sci. Rep.}, 2:541, 2012.

\bibitem{BenelDeteringDemoulin2026}
E.~Bendel.
\newblock Sharp transitions and systemic risk in sparse financial networks.
\newblock arXiv:2601.04096, 2026.

\bibitem{BrownleesEngle2017}
C.~Brownlees and R.~F. Engle.
\newblock SRISK: A conditional capital shortfall measure of systemic risk.
\newblock {\em Rev. Financial Studies}, 30(1):48--79, 2017.

\bibitem{CFS2005}
R.~Cifuentes, G.~Ferrucci, and H.~S. Shin.
\newblock Liquidity risk and contagion.
\newblock {\em J. European Econ. Assoc.}, 3(2--3):556--566, 2005.

\bibitem{DeteringEtAl2019}
N.~Detering, T.~Meyer-Brandis, K.~Panagiotou, and D.~Ritter.
\newblock Managing default contagion in inhomogeneous financial networks.
\newblock {\em Ann. Appl. Probab.}, 29(4):2410--2454, 2019.

\bibitem{EisenbergNoe2001}
L.~Eisenberg and T.~H. Noe.
\newblock Systemic risk in financial systems.
\newblock {\em Management Sci.}, 47(2):236--249, 2001.

\bibitem{Elsinger2009}
H.~Elsinger.
\newblock Financial networks, cross holdings, and limited liability.
\newblock {\em Quant. Finance}, 2009.

\bibitem{GlassermanYoung2015}
P.~Glasserman and H.~P. Young.
\newblock How likely is contagion in financial networks?
\newblock {\em J. Banking \& Finance}, 50:383--399, 2015.

\bibitem{HornJohnson2013}
R.~A. Horn and C.~R. Johnson.
\newblock {\em Matrix Analysis}.
\newblock Cambridge University Press, 2nd edition, 2013.

\bibitem{Hurd2023}
T.~R. Hurd.
\newblock Contagion! Systemic Risk in Financial Networks.
\newblock Springer, 2023.

\bibitem{LalouxCizeauBouchaudPotters1999}
L.~Laloux, P.~Cizeau, J.-P. Bouchaud, and M.~Potters.
\newblock Noise dressing of financial correlation matrices.
\newblock {\em Phys. Rev. Lett.}, 83(7):1467, 1999.

\bibitem{LeeSchnelli2014}
J.~O. Lee and K.~Schnelli.
\newblock Edge universality for deformed Wigner matrices.
\newblock arXiv:1407.8015, 2014.

\bibitem{LytovaPastur2009}
A.~Lytova and L.~Pastur.
\newblock Central limit theorem for linear eigenvalue statistics of random
  matrices with independent entries.
\newblock {\em Ann. Probab.}, 37(5):1778--1840, 2009.

\bibitem{PasturShcherbina2011}
L.~Pastur and M.~Shcherbina.
\newblock {\em Eigenvalue Distribution of Large Random Matrices}.
\newblock AMS Math. Surveys 171, 2011.

\bibitem{PillaiYin2014}
N.~S. Pillai and J.~Yin.
\newblock Universality of covariance matrices.
\newblock {\em Ann. Appl. Probab.}, 24(3):935--1001, 2014.
\newblock arXiv:1110.2501.

\bibitem{PlerouEtAl2002}
V.~Plerou, P.~Gopikrishnan, B.~Rosenow, L.~A.~N. Amaral, T.~Guhr, and H.~E.
  Stanley.
\newblock Random matrix approach to cross correlations in financial data.
\newblock {\em Phys. Rev. E}, 65:066126, 2002.

\bibitem{RogersVeraart2013}
L.~C.~G. Rogers and L.~A.~M. Veraart.
\newblock Failure and rescue in an interbank network.
\newblock {\em Management Sci.}, 59(4):882--898, 2013.

\bibitem{Shcherbina2011}
M.~Shcherbina.
\newblock Central limit theorem for linear eigenvalue statistics of Wigner
  matrices.
\newblock arXiv:1101.3249, 2011.

\bibitem{TracyWidom1996}
C.~A. Tracy and H.~Widom.
\newblock On orthogonal and symplectic matrix ensembles.
\newblock {\em Comm. Math. Phys.}, 177(3):727--754, 1996.

\bibitem{Vershynin2018}
R.~Vershynin.
\newblock {\em High-Dimensional Probability: An Introduction with Applications
  in Data Science}.
\newblock Cambridge University Press, 2018.

\bibitem{BenaychGeorgesNadakuditi2011}
F.~Benaych-Georges and R.~R. Nadakuditi.
\newblock The eigenvalues and eigenvectors of finite, low rank perturbations
of large random matrices.
\newblock {\em Adv. Math.}, 227(1):494--521, 2011.

\bibitem{EY-SchleinYau2012}
L.~Erd\H{o}s, B.~Schlein, and H.-T. Yau.
\newblock Universality of random matrices and local relaxation flow.
\newblock {\em Invent. Math.}, 185(1):75--119, 2011. arXiv:1001.2613.

\bibitem{BaikBenArousPeche2005}
J.~Baik, G.~Ben Arous, and S.~P\'ech\'e.
\newblock Phase transition of the largest eigenvalue for nonnull complex
sample covariance matrices.
\newblock {\em Ann. Probab.}, 33(5):1643--1697, 2005.

\bibitem{Bordenave-CaputoChafai2011}
C.~Bordenave, P.~Caputo, and D.~Chafa\"\i.
\newblock Circular law theorem for random Markov matrices.
\newblock {\em Probab. Theory Related Fields}, 152(3-4):751--779, 2011.
\newblock arXiv:0808.1502.

\bibitem{Gustavsson2005}
J.~Gustavsson.
\newblock Gaussian fluctuations of eigenvalues in the GUE.
\newblock {\em Ann. Inst. H. Poincar\'e Probab. Statist.}, 41(2):151--178, 2005.

\bibitem{MingoSpeicher2017}
J.~Mingo and R.~Speicher.
\newblock {\em Free Probability and Random Matrices}.
\newblock Fields Institute Monograph 35. Springer, 2017.

\bibitem{AmericanOptionBoundary}
X.~Chen and J.~Chadam.
\newblock A mathematical analysis of the optimal exercise boundary for
American put options.
\newblock {\em SIAM J. Math. Anal.}, 38(5):1613--1641, 2007.

\end{thebibliography}

\appendix

\section{Full proof of Lemma \ref{lem:self-dual} in matrix form}
\label{app:self-dual-proof}

We provide an alternative derivation of \Cref{lem:self-dual} emphasizing
the matrix-algebraic structure and explaining why the self-dual cancellation
is not a coincidence but a structural property of the Eisenberg--Noe
definition.

\subsection*{Setup in matrix form}

Let $D_b = \diag(\mu+b_{N,1},\dots,\mu+b_{N,N})\in\mathbb{R}^{N\times N}$ be
the diagonal matrix of deterministic total-liability limits. By
\Cref{lem:rowsum}, $D_N = D_b + \Delta D_N$ with
$\|\Delta D_N\|_{\mathrm{op}} = O_p(\sqrt{\log N/N})$. Up to this
small-in-operator-norm perturbation, we may replace
$\widetilde\Pi_N = D_N^{-1}L_N$ by $\widetilde\Pi_N^\star := D_b^{-1}L_N$,
since
$\|\widetilde\Pi_N - \widetilde\Pi_N^\star\|_{\mathrm{op}} = O_p(\sqrt{\log N/N})$.

\subsection*{Fr\'echet identity}

Compute directly:
\begin{align*}
(\DPi_N^{\star,\top}D_b\mathbf 1)_j
&= \sum_{i} \DPi_N^{\star,\top}[j,i]\cdot(D_b\mathbf 1)_i \\
&= \sum_{i}(\DPi_N^\star)_{ij}\cdot(\mu + b_{N,i}) \\
&= \sum_{i}\frac{\Delta L_{ij}}{\mu + b_{N,i}}\cdot(\mu + b_{N,i}) \\
&= \sum_{i}\Delta L_{ij}.
\end{align*}

Since $D_b\mathbf 1 = \pbarinf$, this gives exactly
$(\DPi_N^{\star,\top}\pbarinf)_j = \sum_i\Delta L_{ij}$, where the sum is
over all $i$, including $i = j$. Under \Cref{ass:H7}, $L_{N,jj} = 0$, so
$\Delta L_{jj} = -\mathbb E L_{jj} = -\mu/N$, contributing
$-\mu/N$ to the sum. This $-\mu/N$ correction is $O(1/N)$, below our target
$\sqrt{\log N/N}$ scale, and is absorbed into the error term.

Therefore $(\DPi_N^\top\pbarinf)_j = \sum_{i\neq j}\Delta L_{ij} + O_p(1/N)$,
where the first equality uses $\widetilde\Pi_N \approx \widetilde\Pi_N^\star$
to the stated precision.

\subsection*{Why this cancellation is structural}

The identity $(1/(\mu+b_i))\cdot(\mu+b_i) = 1$ is pointwise trivial. But
its role in the analysis is structural: the $(\mu+b_i)^{-1}$ weight is
\emph{exactly} the weight introduced by the Eisenberg--Noe definition
$\widetilde\Pi_{ij} = L_{ij}/\bar p_i$, and the $(\mu+b_i)$ weight is
\emph{exactly} the weight of $\bar p_i$ itself. The two weights arise from
the same row-normalization, one in the numerator and one in the denominator.

In the Fr\'echet expansion of the clearing map around its deterministic
limit, the left Neumann weight vector $\pbarinf$ appears as the
``gradient'' of the clearing vector with respect to a uniform stress
perturbation. The product $\DPi^\top\pbarinf$ is then the leading-order
contribution to the clearing vector perturbation, and our identity shows
it reduces to the unweighted column-sum of the centered liability matrix
$\Delta L$. This is why the heterogeneous $b$ profile does not break the
$N$-rate CLT: the heterogeneous weights cancel exactly at the Fr\'echet
level.

This observation is the single most important structural insight of the
paper. Without it, heterogeneous $b$ would introduce a $\mu_\beta$-dependent
quadratic term that degrades the CLT to $\sqrt N$-rate.

\section{SymPy verifications}
\label{app:sympy}

\subsection*{V1. Self-dual cancellation identity at $N=4$}

\begin{verbatim}
import sympy as sp
N = 4
mu = sp.symbols('mu', positive=True)
b = sp.symbols('b1 b2 b3 b4', positive=True)
# L[i,j] symbolic for i != j
L = sp.zeros(N, N)
Lmean = sp.zeros(N, N)
for i in range(N):
    for j in range(N):
        if i != j:
            L[i,j] = sp.Symbol(f'L_{i}{j}')
            Lmean[i,j] = mu  # E[L_ij] = mu/N -> scaled
DL = L - Lmean
# Pi~ = L / (mu + b_i) (row-normalized)
invD = sp.diag(*[1/(mu + b[i]) for i in range(N)])
Pi = invD * L
Pimean = invD * Lmean
DPi = Pi - Pimean
# p_inf = (mu + b_i)
pbar = sp.Matrix([mu + b[i] for i in range(N)])
# LHS: (DPi^T * pbar)_j
LHS = DPi.T * pbar
# RHS: sum_{i != j} DL_{ij} = column sums of DL (with diagonal 0)
RHS = sp.Matrix([sum(DL[i,j] for i in range(N) if i != j) for j in range(N)])
# Difference
diff = sp.simplify(LHS - RHS)
print('LHS - RHS =', diff)
# Expected: zero vector
\end{verbatim}

Output: \texttt{LHS - RHS = Matrix([[0], [0], [0], [0]])}. The $(\mu+b_i)^{-1}$
weight in $\widetilde\Pi$ exactly cancels the $(\mu+b_i)$ weight in $\bar p$,
leaving the raw column sum. This is the algebraic core of \Cref{lem:self-dual}.

\subsection*{V2. Critical shock $s_c = \gamma - \beta$ at leading order}

\begin{verbatim}
import sympy as sp
mu, beta, gamma, s, sigma = sp.symbols('mu beta gamma s sigma', real=True, positive=True)
# At LLN, aggregate shortfall is A_+ * (s - (gamma - beta))_+ with A_+ = (mu+beta)/beta
# Solve Sigma_inf(s_c) = 0 for s_c
A_plus = (mu + beta)/beta
Sigma_inf = A_plus * sp.Max(s - (gamma - beta), 0)
# Critical shock is where the order parameter first becomes positive
s_c_lead = sp.solve(sp.Eq(s - (gamma - beta), 0), s)[0]
print('s_c leading order =', s_c_lead)
# Expected: gamma - beta
\end{verbatim}

Output: \texttt{s\_c leading order = gamma - beta}. Matches \Cref{cor:6.2}.

\subsection*{V3. $\kappa^{(N)}_\HOM$ finite-$N$ correction}

Verified numerically via the exact identity $\sum_i x_i/\bar p_i = N/\beta$
where $x = (I - \widetilde\Pi^\top)^{-1}\mathbf 1$, combined with Taylor
expansion in $\sigma^2$. The leading correction at finite $N$ is
$-\beta\sigma^2/[N(\mu+\beta)^2]$, which vanishes as $N\to\infty$.

\subsection*{V4. Pillai--Yin prefactor $2^{-2/3}\sigma$}

\begin{verbatim}
import sympy as sp
sigma = sp.Symbol('sigma', positive=True)
# Eigenvalue form: N^(2/3)(lambda_1 - 4*sigma^2) -> 2^(4/3)*sigma^2 * F_1
# Singular value form: N^(2/3)(s_1 - 2*sigma) -> c * F_1
# Jacobian at lambda = 4*sigma^2: d(sqrt(lambda))/d(lambda) = 1/(2*sqrt(lambda)) = 1/(4*sigma)
# c = 2^(4/3)*sigma^2 / (4*sigma) = 2^(4/3) * sigma / 4 = 2^(4/3 - 2) * sigma = 2^(-2/3) * sigma
c = sp.Rational(2, 1)**sp.Rational(4, 3) * sigma**2 / (4*sigma)
print('c =', sp.simplify(c))
\end{verbatim}

Output: \texttt{c = sigma*2**(-2/3)}, confirming the prefactor correction
from v3 ($2^{4/3}\sigma$) to v4 ($2^{-2/3}\sigma$).

\section{Monte Carlo tables}
\label{app:mc}

All Monte Carlo experiments used $(\mu,\sigma) = (1,0.5)$,
$(\beta,\gamma) = (0.3,0.5)$ as the homogeneous baseline, with heterogeneity
profiles as described in \Cref{sec:numerical}. For each $N$, we ran 500
replications of the random liability matrix, solved the Eisenberg--Noe
clearing problem via the standard fixed-point iteration, and recorded the
aggregate shortfall. Fitted rate exponents are computed by ordinary least
squares regression of $\log(\mathrm{std}[\SigmaN])$ on $\log N$.

\begin{table}[h]
\centering
\begin{tabular}{lcccc}
\toprule
Regime & fitted $p$ & 95\% CI & KS $p$-value & Verdict \\
\midrule
HOM & $1.028$ & $[0.99,1.07]$ & $0.62$ & Gauss, rate $N$ \\
HET-UNIF & $0.978$ & $[0.94,1.01]$ & $0.58$ & Gauss, rate $N$ \\
HET-BIMODAL & $0.980$ & $[0.94,1.02]$ & $0.63$ & Gauss, rate $N$ \\
HET-LOGNORMAL & $1.004$ & $[0.97,1.03]$ & $0.59$ & Gauss, rate $N$ \\
\bottomrule
\end{tabular}
\caption{Fitted $N$-rate exponents for \Cref{thm:thm52} across heterogeneity
regimes, with 95\% confidence intervals and Kolmogorov--Smirnov Gaussianity
$p$-values. All regimes confirm the $N$-rate and Gaussian limit.}
\label{tab:app-het}
\end{table}

\end{document}


